% Generated by report_gen.py (ai-ee v1) - do not hand-edit.
\documentclass[11pt,a4paper]{article}
\usepackage[margin=2.2cm]{geometry}
\usepackage{graphicx}
\usepackage{booktabs}
\usepackage{longtable}
\usepackage{pdfpages}
\makeatletter
\@ifundefined{KV@Gin@artifact}{\define@key{Gin}{artifact}[]{}}{}
\makeatother
\usepackage{xcolor}
\usepackage[hidelinks]{hyperref}
\setcounter{tocdepth}{1}
\begin{document}
\sloppy
\begin{center}
{\LARGE\bfseries bb-mcu --- Design Document}\\[6pt]
{\large ai-ee v1 pipeline}\\[2pt]
generated 2026-08-17 20:50:41
\end{center}
\tableofcontents

\section{Board and Run Metadata}
{\small
\begin{longtable}{lp{11cm}}
\toprule
\textbf{Field} & \textbf{Value} \\
\midrule
\endhead
board & bb-mcu \\
workspace & \texttt{boards/bb-mcu} \\
phase & P9 \\
gate status & all 5 recorded gates pass \\
state created & 2026-08-16T17:25:37 \\
state updated & 2026-08-17T20:49:20 \\
\bottomrule
\end{longtable}
}
\section{Overview}
learning block-basics: a microcontroller board. Powered from an external 3.3 V rail on a screw terminal - no regulation on this board. Programmable and debuggable over the MCU's standard debug interface, with four GPIO brought out to a 0.1 inch header. Any MCU sourceable at JLCPCB is acceptable. Workspace boards/bb-mcu.

---

Owner's brief, verbatim, received 2026-08-16.

\section{Requirements}
\subsection*{Requirements: bb-mcu}
\subsubsection*{1. Function}
A bare microcontroller board, built as a study article. It takes a ready-made 3.3 V rail from an external supply on a screw terminal - there is no regulation, no conversion and no second rail on this board - and gives that MCU three things and nothing else: power in, the MCU's own standard debug/programming interface brought out to a header, and four general-purpose I/O brought out to a 0.1 inch header. The board's whole job is to hold one MCU, let it be programmed and debugged, and let four of its pins be reached.

BUILD MODE (declared by the brief's opening token, resolved and recorded as the P0 decision in \texttt{state.json}; contract in \texttt{reference/build-modes.md}):

\begin{itemize}
\item token: \texttt{learning block-basics:}
\item target learning outcome: \textbf{block-basics} - the block end to end
\item scope tier: \textbf{block-only}
\item binding level: \textbf{canonical}
\item stage under study: \textbf{none} (no single stage is being taught; the whole
P0-P9 run is the lesson)
\item geometry: \textbf{OUTPUT} - the brief states no size, none is invented here, and
the placement earns the outline at P6 (see section 5)
\end{itemize}
The block under study is the MCU itself. In scope: the MCU, plus exactly the support components its datasheet requires for correct operation at the stated operating point (decoupling, and whatever the part needs on NRST / boot-mode / supply pins to run and to be debugged), plus the interfaces named below, plus what the fab needs to build the board and the bench needs to hold it. Out by mode - a SCOPE decision, not an engineering conclusion, and never a reviewer finding: protection of every kind, filtering the datasheet does not require, indicators (no LEDs), buttons (no reset button), test points, config straps and jumpers, any second rail or second IC, and mechanical/enclosure features beyond mounting the bare board. The mode's defaults are APPLIED below (quantity 5, JLC PCBA single-sided, indoor bench 0-50 C, no enclosure, fewest honest layers, size earned by the layout, stop at P9) instead of being asked.

On "one input interface, one output interface": the input interface is the 3.3 V power in on the screw terminal, and the output interface is the 4-way GPIO header. The debug interface is the third connector and it STAYS - it is not a test point and not scope creep. An MCU that cannot be programmed is not a working MCU, so its standard debug/programming access is part of the block under study, and the brief names it explicitly. All three connectors below are stated by the brief; nothing beyond them is added.

ASSUMED: this pipeline delivers hardware only. No firmware, no bootloader image and no application program is a deliverable here; "programmable" is a property of the board, not a software requirement on the run.

\subsubsection*{2. Interfaces}
Three, all stated by the brief.

\begin{itemize}
\item \textbf{J1 - power input, screw terminal (STATED).} 2-pin, DC only, 3.3 V from
an external rail. ASSUMED: through-hole 2-pole terminal block, 5.08 mm
pitch, the well-stocked JLC/LCSC choice that takes bench lead wire (pitch
was not stated; packaging is relaxable under this mode, so it is assumed
rather than asked). Polarity silkscreen-marked (+ / -). NOTE: there is no
reverse-polarity protection by mode, so the silk marking is the only
defense against a swapped supply.
\item {*}{*}J2 - debug / programming header (STATED as "the MCU's standard debug
interface").{*}{*} ASSUMED: 0.1 inch pin header, single row, straight,
through-hole, carrying whatever the chosen MCU's standard interface is -
for a Cortex-M part that is SWD (SWDIO + SWCLK), with GND, and with NRST
brought out if the part's standard programming flow expects it. ASSUMED: no
trace/SWO pin and no UART bootloader header - the standard debug interface
alone is what the brief asked for. Whether the header also carries 3.3 V as
a probe reference pin is open question 2. ASSUMED: any common flying-lead
SWD probe (ST-LINK / J-Link / DAPLink class) is the tool; no vendor-specific
10-pin or 20-pin box connector, since packaging is relaxable here.
\item {*}{*}J3 - GPIO header (STATED: "four GPIO brought out to a 0.1 inch
header").{*}{*} 0.1 inch pitch, single row, through-hole. ASSUMED: 5 positions
\begin{itemize}
\item the 4 GPIO plus one GND - because a signal with no return is not usable
at the bench; this is the return path of the stated interface, not an added
feature. Electrical expectations for the four pins are open question 3.
\end{itemize}
\item \textbf{No other external interface exists or is wanted}: no USB, no comms
connector, no LED, no button, no jumper, no test point, no external
reset, no expansion header. All excluded by mode.
\end{itemize}
\textbf{What the MCU choice must satisfy} (the brief says "any MCU sourceable at JLCPCB is acceptable", so the part number is a P1/P3 selection job, not an owner question - these are the constraints the selection is bound by, and they are recorded here so the selection is checkable):

1. Single supply, runs correctly across at least 3.0-3.6 V (see section 3); no internal or external regulator needed on this board. 2. A standard hardware debug/programming interface on dedicated pins, usable with a common third-party probe, and NOT one that has to be multiplexed away from the four GPIO. 3. At least four free general-purpose I/O left over after the debug pins and any pins the datasheet reserves. 4. Runs and debugs on its internal oscillator - ASSUMED no external crystal, because nothing in the brief needs timing accuracy. A part that mandates an external clock to be programmed is disqualified rather than accompanied by a crystal. 5. No mandatory external components beyond decoupling and whatever the datasheet requires on NRST / boot-mode pins. IMPORTANT: if the chosen part needs a boot-mode pin held at a defined level to run or to be debugged, that resistor is datasheet-required SUPPORT and is in scope - it is not an excluded "config strap". 6. Sourceable and assemblable at JLCPCB: in stock, JLC Basic or Preferred where one meets the above, and in a package their PCBA line places on the top side.

\subsubsection*{3. Power}
\begin{itemize}
\item \textbf{Source: an external 3.3 V rail on J1 (STATED).} What generates that rail
is NOT stated by the brief - see section 8 and open question 1. There is no
regulation, conversion, sequencing or power switching on this board
(STATED), and no second rail (also excluded by mode).
\item ASSUMED: the rail is 3.3 V nominal at +/-5 \% (3.14-3.47 V), and the part
selection must tolerate at least 3.0-3.6 V so an ordinary "3.3 V" bench or
dev-board rail is always inside range. Confirmed or corrected by the answer
to question 1.
\item \textbf{Rail budget (GUESS, for sizing only):} a small MCU on its internal clock
draws roughly 5-30 mA active at 3.3 V, under \textasciitilde{}100 mA even for a large part
running flat out with peripherals on. Board total GUESSED at \textless{} 100 mA, i.e.
\textless{} 0.35 W. Nothing on this board is a thermal problem, and the terminal, the
headers and the copper are all trivially oversized for this current at any
sane geometry.
\item Decoupling per the chosen MCU's datasheet (per-supply-pin ceramics plus any
bulk it specifies) is datasheet-required support and is IN scope. Filtering
the datasheet does not require - ferrite beads, pi filters, a separate
analog supply network the part does not ask for - is excluded by mode.
\item No battery on this board, no charging function, no energy storage, no backup
rail, no supercap, no RTC coin cell. None is stated and none is implied by
"an external 3.3 V rail".
\item Back-feed note (not a safety flag, a bench fact): if question 2 is answered
"yes, put 3.3 V on the debug header", then the board has two ways in and no
ORing or protection by mode. Powering it from the probe and the screw
terminal at the same time is prevented by bench procedure only - the board
will not enforce it.
\item Excluded by mode (SCOPE decisions - do not record any of these as generally
unnecessary): input TVS/ESD, reverse-polarity protection, fuse, UVLO/OVP,
input EMI filter, power-good indicator, supply sequencing.
\end{itemize}
\subsubsection*{4. Environment}
Mode defaults - applied, not asked.

\begin{itemize}
\item Ambient: 0 to 50 C, indoor bench.
\item Cooling: natural convection. No forced airflow, no heatsink, no chassis
thermal path. Nothing on this board dissipates enough to need one.
\item No enclosure. No ingress (IP) rating. No vibration, shock, humidity or
conformal-coating requirement.
\item No formal EMC/emissions campaign - this is a bench study article, not a
product. Sane decoupling and return-path layout are still required because
they are what make the block correct, not because a test demands them.
\end{itemize}
\subsubsection*{5. Size \& mounting}
\begin{itemize}
\item \textbf{Outline: no HARD cap, and none may be introduced later.} The brief
states no dimension, and none is invented here. Under binding \texttt{canonical}
the board size, aspect and outline are OUTPUTS of the placement:
\textbf{RELAXABLE (canonical)} applies to every geometric preference on this
board, and any dimension that appears at a later checkpoint is a preference
that LOSES to the canonical layout, with the loss recorded as a
\texttt{state.py decision}.
\item Flow that follows from that binding: \texttt{board\_init --outline auto} at P5
(generous provisional room - it must NOT be given a fixed size; it will
refuse one), place to the canonical layout at P6 and pass the \texttt{place} gate,
then \texttt{board\_edit --outline fit --margin M} so the board becomes what the
placement needs (re-run \texttt{planes\_gen} if it grew), then route at P7.
\item The honest layout here is driven by the three through-hole connectors and
by keeping the MCU's decoupling tight to its supply pins - not by heat and
not by any stated dimension. No planning number is offered deliberately:
on this board a number would be an anchor with nothing behind it.
\item Layers: ASSUMED 2, per the mode's "fewest layers the block honestly needs".
An MCU with decoupling, three connectors and a ground pour is a genuine
2-layer board. Not fixed here - if the P2/P4 work shows otherwise, 4 layers
is allowed under the same rule.
\item Mounting: ASSUMED four M3 clearance holes (3.2 mm) inset from the corners
so the board can sit on standoffs at the bench; two on opposite corners is
acceptable if the earned outline is too small for four. Mounting the bare
board is in mode scope; anything beyond it (standoffs, brackets, enclosure
bosses) is not.
\item Height: ASSUMED no maximum. The board is open on a bench with no enclosure,
and the screw terminal itself stands \textasciitilde{}10-12 mm with wire fitted.
\item Connector placement: ASSUMED all three connectors on board edges with their
openings facing outward, power on a different edge from the two signal
headers, so bench wiring and the probe do not cross the board. Which edges
is the layout engineer's call at P6.
\end{itemize}
\subsubsection*{6. Quantity \& budget}
Mode defaults - applied, not asked.

\begin{itemize}
\item Build quantity: 5.
\item Cost: minimal, no hard per-unit cap. Prefer JLC Basic/Preferred parts
wherever one actually meets the requirement; Extended is allowed for the
MCU if that is where the honest choice lives.
\end{itemize}
\subsubsection*{7. Assembly}
Mode defaults - applied, not asked.

\begin{itemize}
\item JLCPCB PCBA, single-sided: all SMT parts on the top side only. The bottom
side stays clear of SMT so it can carry unbroken return copper.
\item The screw terminal and both 0.1 inch headers are through-hole. ASSUMED:
fitted by JLCPCB through-hole assembly if offered at acceptable cost for
the chosen parts, hand-soldered after SMT otherwise. Either path leaves the
schematic and the footprints unchanged.
\item Pipeline stops at P9 (fab package + DFM). Ordering is a separate owner
decision and is not requested here.
\end{itemize}
\subsubsection*{8. Compliance/safety flags}
The mode grants no silence in this section. One flag is genuinely OPEN and must be answered before the pipeline proceeds; the rest are closed by facts the brief states.

\begin{itemize}
\item {*}{*}Batteries: CLOSED 2026-08-16 by owner answer to question 1 - NOT a
battery{*}{*} (bench supply or another board's 3.3 V rail). The reasoning that
made it a question is kept below because it is what the answer closed.
The brief names the
supply as "an external 3.3 V rail" but never says what generates it. A
bench supply, a dev board's 3.3 V pin and a battery pack all present that
way, and battery chemistry/charging is the one thing this pipeline will not
proceed on as a guess. There is no charging circuit and no cell on this
board under any answer; the answer decides whether the source has to be
treated as a battery (discharge behaviour, current limit, deep-discharge)
in later phases. Answer it and this flag closes.
\item \textbf{Mains voltage: DOES NOT APPLY.} The board sees 3.3 V DC only and has no
connection to anything mains-referenced. Generating the 3.3 V is off this
board and out of scope; if the answer to question 1 is "an off-line
supply", the isolation lives in that supply, not here - but say so, because
it changes nothing on this board and everything about how it is handled.
\item \textbf{High current (\textgreater{}3 A): DOES NOT APPLY.} Maximum current anywhere on the
board is the MCU's own supply current, GUESSED under 100 mA (section 3),
more than an order of magnitude below the threshold. Note for question 1:
the SOURCE may be able to deliver far more than the board draws, and this
board has no fuse or current limit by mode. At 3.3 V that is not a hazard
to a person; it is a bench fact worth stating, so it is stated.
\item \textbf{\textgreater{}30 V: DOES NOT APPLY.} Nothing on the board exceeds 3.3 V. Not a
boundary case. \texttt{check\_creepage} stays a clean no-op.
\item \textbf{Motors: DOES NOT APPLY as specified.} Nothing on this board drives a
load. If the owner intends one of the four GPIO to drive a motor, a relay
or any other inductive load (see question 3), this flag re-opens - a GPIO
driving an inductive load off-board needs a conversation about what comes
back up that wire, and this board has no protection by mode.
\item \textbf{RF transmit: DOES NOT APPLY.} No radio function, no antenna, no
transmitter. An MCU with an unused integrated radio would not change this,
but a wireless MCU should not be chosen for that reason alone.
\end{itemize}
\subsubsection*{9. Open questions}
Three. Each has a default that is a sensible answer - reply "defaults" and the pipeline proceeds on all three as written.

1. \textbf{What actually supplies the 3.3 V?} (This is the safety question - it is the one thing here that will not be guessed.) A bench power supply set to 3.3 V, the 3.3 V pin of another powered board, a USB-powered breakout, or a battery pack? If it is a battery of any kind, say which chemistry and how many cells, and whether anything charges it - nothing on this board would. DEFAULT: a bench supply or another board's 3.3 V rail, indoors, current limited to roughly 0.5 A or less, and NOT a battery. On that default the rail is taken as 3.3 V +/-5 \% and the MCU is chosen to accept at least 3.0-3.6 V. 2. \textbf{Should the debug header carry the board's 3.3 V as a reference pin?} Most debug probes want to sense the target's supply voltage before they talk to it, and that sense pin is normally wired to the target's own rail. Including it costs one header position and creates a second path to the power rail on a board that has no protection: if the probe is one that can OUTPUT power on that pin, having the screw terminal connected at the same time is a bench-procedure rule the board cannot enforce. DEFAULT: yes - one 3.3 V pin on the debug header, wired to the same rail, for the probe to SENSE only. The board is always powered from the screw terminal, never from the probe. Say no if you would rather have the header carry only the debug signals and ground. 3. {*}{*}What are the four GPIO for - do any of them need a particular capability?{*}{*} Whether they are plain on/off pins or need to be, say, a UART transmit/receive pair, a PWM-capable output, or an analog input decides WHICH pins they are on the chip, and can decide which chip. The same question covers what gets connected to them: a scope probe or a logic analyser needs nothing extra, but driving an LED, a relay or a motor from one of these pins changes section 8 (see the motors flag). DEFAULT: four plain digital I/O, 3.3 V logic, no alternate-function requirement, no 5 V tolerance, each loaded no harder than the MCU's own per-pin current limit (a few mA - a scope probe, a logic analyser, or an LED with its resistor sitting off-board). Nothing inductive is driven.

\paragraph*{ANSWERS (owner, 2026-08-16) - all three defaults taken}
All three questions are CLOSED. Each answer is also a \texttt{state.py decision} in \texttt{state.json}; this section is the design-facing copy. Nothing below is an assumption any more - later phases treat these as stated requirements.

1. {*}{*}3.3 V source: a bench supply or another board's 3.3 V rail. NOT a battery.{*}{*} Indoors, current-limited to roughly 0.5 A or less; nothing charges and nothing stores energy on this board. The rail is 3.3 V +/-5 \% and the MCU must accept at least 3.0-3.6 V. Section 8's battery flag is closed; the motors and mains flags stay closed too. 2. \textbf{Debug header carries 3.3 V, SENSE-only.} J2 is a 5-position 0.1 inch single-row header carrying \texttt{3V3}, \texttt{SWDIO}, \texttt{SWCLK}, \texttt{NRST}, \texttt{GND}. The owner fixed the SET; the ORDER was left to the layout engineer and was RULED at P1 as \textbf{`GND / SWCLK / 3V3 / SWDIO / NRST`} (position 1 to 5) - see the P1 decision in \texttt{state.json}. 3V3 sits at the centre because that is the only arrangement in which a 180 degree reversal of a 5-way shell lands the probe's high-Z VTref input on the board's 3.3 V rail instead of a probe OUTPUT; GND and NRST take the ends so a reversed probe ground clamps reset rather than a data line. The pinout must still be silk-labelled per pin - the header has no keying, and the labels are what a hand-wired flying lead is read from. The probe senses the rail to set its I/O level. The board is powered from J1 only, never from the probe - a bench-procedure rule the board does not enforce, as section 3's back-feed note says. 3. \textbf{The four GPIO are plain digital I/O.} 3.3 V logic, no alternate-function requirement, no 5 V tolerance, no inductive load, each loaded no harder than the MCU's own per-pin limit. Pin selection at P4 is therefore free to pick whichever four ordinary GPIO give the cleanest escape from the package - that is now a LAYOUT criterion, not an electrical one.

\paragraph*{DELEGATION (owner, 2026-08-16)}
"Make decisions yourself based on what will give the best learning for the basic system." Design judgment calls in this run are ruled by the orchestrator against that criterion and recorded as decisions; H1-H4 are presented as digests rather than blocking approvals. This delegation relaxes NOTHING that makes the board true: every gate, the P2/P3 coverage checks and the research they trigger, DFM, the datasheet's own requirements and every section 8 safety question stand exactly as written. Money and irreversible steps still go to the owner - and this board stops at P9 by mode anyway.

\section{Architecture}
\subsection*{Decisions of record (state.json)}
{\small
\begin{longtable}{lp{5.2cm}p{7.2cm}}
\toprule
\textbf{Phase} & \textbf{What} & \textbf{Why} \\
\midrule
\endhead
P0 & Build mode: mode block-basics; scope block-only; binding canonical & token 'learning block-basics:' in the brief; contract in reference/build-modes.md \\
P0 & 3.3 V source is a bench supply or another board's 3.3 V rail - NOT a battery; indoors, current-limited \textasciitilde{}0.5 A; nothing charges on this board & owner answer to P0 safety question 1; closes the section-8 battery flag. Rail spec: 3.3 V +/-5 percent, MCU must accept 3.0-3.6 V \\
P0 & Debug header is a 5-pin 0.1 inch flying-lead pinout: 3V3 / SWDIO / SWCLK / NRST / GND, the 3V3 pin SENSE-only & owner answer to P0 question 2; probe reads the rail to set its I/O level, board is always powered from the screw terminal and never from the probe \\
P0 & The four GPIO are plain digital I/O: 3.3 V logic, no alternate-function requirement, no 5 V tolerance, nothing inductive driven & owner answer to P0 question 3; keeps MCU pin selection unconstrained and keeps the motors/inductive-load safety flag closed \\
P0 & Owner delegated design rulings to the orchestrator for this run: 'Make decisions yourself based on what will give the best learning for the basic system' & H1-H4 are presented as digests and ruled by the orchestrator with that criterion; owner is still consulted for safety questions and any money/irreversible step. Gates, coverage, research and DFM are unaffected - a mode and a delegation relax neither \\
P1 & P1 roster: research-component-scout (wave A, alone), then research-reference-design + research-interface-spec in parallel (wave B) off the scout's shortlist. research-power-architect SKIPPED & the MCU family decides what the 'standard debug interface' even is (ARM SWD vs a RISC-V single-wire port), so interface-spec and reference-design must not run blind - wave A gates them. Power-architect is skipped under its own 'unless trivially powered' clause: one external 3.3 V rail, no conversion, no sequencing, no second rail, \textless{}100 mA guessed \\
P1 & MCU family ruled: STM32F030F4P6TR (LCSC C89040, TSSOP-20, ARM Cortex-M0, SWD on dedicated PA13/PA14). Runners-up rejected: STM32C011F4P6, PY32F030F28P6TU, CH32V003F4P6, STM32G031F4P6 & learning criterion (owner delegation). SWD is the canonical standard debug interface the brief names; CH32V003 single-wire SWIO needs a WCH-Link-class probe and would rewrite the owner-answered J2 pinout. PY32 is disqualified on a MECHANICAL ground for a library-seeding board: puya.com is not on the research fetch allowlist, so its records could never be sourced - st.com is. Against C011: F030 has genuinely dedicated NRST and BOOT0 pins (no alt function on either), so the minimum system teaches boot-mode strapping as its own concept instead of hiding it on SWCLK, and its required 10k BOOT0 pull-down is the exact 'datasheet-required support, not an excluded config strap' boundary this scope tier turns on. Stock 19015 vs 1371 and the deepest documentation base of any 20-pin MCU seal it. Price is not a driver at qty 5 \\
P1 & J2 pin order RULED: GND / SWCLK / 3V3 / SWDIO / NRST (position 1..5). Overrides both the order listed in requirements.md section 9 answer 2 and the interface-spec agent's recommended 3V3/SWCLK/GND/SWDIO/NRST & the owner fixed the SET of pins and left the ORDER to the layout engineer, so this uses that latitude, it does not contradict the answer. With a 5-way shell a 180 deg reversal maps probe wire i onto board position 6-i, so the only destructive outcome is a probe OUTPUT landing on board 3V3. Board 3V3 receives the probe's high-Z VTref input if and only if 3V3 sits at position 3 - that makes 3V3-at-centre the unique reversal-safe arrangement. The stated order puts probe GND on 3V3 (hard short); the agent's GND-at-centre order puts the probe's open-drain NRST on 3V3 (shorts whenever reset is asserted, which connect-under-reset does every time). GND and NRST take the ends so probe GND clamps reset (benign, obvious at the bench) instead of a data line, and GND stays adjacent to SWCLK per ARM's return-next-to-clock convention. Cost accepted: this is deliberately NOT a vendor cable order - real flying-lead probes present individually labelled wires, so per-pin silk is what the user reads either way \\
P1 & No series resistors on SWDIO/SWCLK. Considered rejection, not an absence of guidance & Raspberry Pi (100R), Lauterbach (47R) and SEGGER all recommend a series resistor on the target-driven line - this ruling overrules them knowingly and must never be defended with 'no source asks'. Grounds: AN4325 4.3.3 states the embedded pull-up/pull-down remove the need for external resistors; the fastest edge the part can produce is 5 ns (datasheet Table 48) and reflections on a \textless{}=50 mm FR-4 trace settle in \textasciitilde{}1.5 ns against a 125 ns half-bit at ST-LINK/V2's 4 MHz ceiling; and a damping/limiting resistor the datasheet does not require is exactly the filtering that scope tier block-only excludes. Keeping the BOM at what ST actually requires is what makes the required-vs-recommended line visible on this board \\
P1 & No 100 nF capacitor on NRST. Diverges from this repo's stm32-blinky board, which fitted one (C9) & AN4325 files it under Table 6 Optional components, 'for RESET button', and this board has no reset button (buttons are excluded by mode). NRST's permanent internal \textasciitilde{}40k pull-up needs no external part and the die already filters pulses \textless{}=100 ns. Accepted risk, stated so bring-up can judge it: NRST is brought out to J2 on a flying lead, so a noise-triggered reset is conceivable - if bring-up sees one, THAT is the finding this board earned, and it is a knowledge record, not a design error to pre-empt \\
P1 & constraints.json OMITS the diff\_pairs key entirely rather than declaring diff\_pairs: {[}{]} & per constraints\_schema.md an explicit empty list DISABLES check\_diffpair board-wide while omitting the key leaves auto-discovery armed. Both are no-ops on a board with no pair, so take the fail-safe one: disabling a check board-wide to say 'we have none' is how a silent blind spot gets built in for the next edit \\
P1 & SWD binds NOTHING on layout: no controlled impedance, no pair, no length matching, no termination, no return-path rule. high\_speed and diff\_pairs stay empty of SWD entries & 5 ns worst-case edge settles in \textasciitilde{}1.5 ns on a \textless{}=50 mm trace against a 125 ns half-bit at 4 MHz (80x margin, still 14x at STLINK-V3's 24 MHz), and stackups.yaml records live-verified that JLC sells NO impedance-controlled 2-layer stackup - an impedance\_ohm key here would be unbuildable as well as unnecessary. The \textless{}=50 mm / over-the-pour / keep-SWCLK-clear guidance stays ADVISORY in notes, never an enforced constraint. Recorded because the temptation to invent a 50 ohm target is exactly what would propagate a fake requirement into P5 rules and P7 routing \\
P2 & constraints.json declares TWO blocks{[}{]} entries - topology 'mcu' and topology 'swd-debug-port' - not one & coverage interface slots are derived ONLY from diff\_pairs{[}{]}.base (knowledgelib.workspace\_slots), so SWD would get no slot at all unless a fake differential pair were invented, which would inject an unbuildable impedance requirement into P5/P7. Declaring the debug port as its own block is the honest way to make the machinery cover it: it IS a functional block of this board, and its knowledge (pull requirements, reversal-safe pinout, header conventions) is reusable across ANY MCU board independent of the MCU family, which is exactly the kind of record this seeding run exists to produce. Cost: one extra research task against a per\_run budget of 6, with 3 expected in total \\
P2 & J3 pin order: IO1 / IO2 / GND / IO3 / IO4 (positions 1..5). Four GPIO = PA0, PA1, PA2, PA3 (TSSOP-20 pins 6-9) & owner fixed the SET (4 GPIO + a GND), order was open. GND at centre halves the worst-case signal-to-return distance on the flying lead (4 positions to 2) and reuses J2's reversal geometry; nothing destructive either way here since J3 carries no rail. PA0-PA3 chosen as a pure LAYOUT criterion (the owner's answer 3 freed it from any electrical one): BOOT0(1)/NRST(4)/SWDIO(19)/SWCLK(20) all cluster at ONE package end, so J2 takes all three debug signals off that end without crossing, while PA0-PA3 escape from the diagonally opposite region. Tie-breakers over the equally adjacent PA5-PB1: one port with consecutive bits, and maximum distance from SWCLK. No alternate-function claim is made - this run's research did not verify the AF map \\
P2 & Stackup JLC2313\_1.6, 2 layers, 1 oz HASL, B.Cu unbroken GND pour by planes\_gen default. No 4-layer escalation condition written. One flat root sheet, hierarchy rejected & verified available:true and the defaults{[}2{]} entry in stackups.yaml. JLC sells no impedance-controlled 2-layer stackup and nothing on this board needs one - that is the layer-count justification, not an apology. A 4L trigger would have nothing to trigger: a TSSOP-20 with no exposed pad has no thermal path into an inner plane to buy. Hierarchy rejected for nine parts and ten nets on the recorded /\textless{}sheet\textgreater{}/\textless{}LABEL\textgreater{} export hazard that cost lumina-par \\
P2 & power{[}{]} declares +3V3 only, NOT GND. placement.sides pins the six SMT parts (U1, C1-C4, R1) to front; the three THT connectors get no entry & declaring GND would manufacture a P8 error: at 0.1 A the declaration sizes nothing (rules\_gen buckets it into Default) but check\_pdn errors pdn\_undecoupled on a rail with no caps, and the only fix would be a waiver. Deliberate divergence from bb-buck, which wrote exactly that waiver at 2.6 A where it was worth it. The sides entries are what enforce requirements s7 single-sided PCBA against the U19 annealer, which can otherwise legitimately move a free SMT cluster to the back; I verified the THT exemption directly at place\_anneal.py:250, flippable=(kind=='free' and not thru), so a through-hole part is structurally unflippable and needs no entry \\
P2 & Connector edges declared as a rotatable SET: J1 right (pos 0.5), J2 top, J3 left. rot deliberately unset on all three. NOTHING was relaxed by the mode at P2 & what binds is three DIFFERENT edges with power apart from the two signal headers, not the compass directions; J1's pos 0.5 is load-bearing (directly opposite VDD/VSS pins 15/16 and clear of the corner mounting holes). rot is left unset because bb-buck recorded the cost of guessing a footprint's native entry direction. No relaxation decision exists because the brief states no dimension - there is no stated value to lose, and geometry stays an OUTPUT earned at P6 \\
P2 & VDDA is treated as this board's CLOCK supply, tied directly to VDD with no series element of any kind. Any ferrite/filter into VDDA is a design FAULT here, not merely an out-of-scope extra & P2 research found datasheet Figure 12 (p41) puts the internal RC oscillators and the PLL in the analog domain fed from VDDA - and this board runs on the internal RC with no crystal, so VDDA powers the clock, not an unused ADC. Table 21 additionally states VDDA must sit at a potential equal to or higher than VDD, which a series element can only violate. This REPLACES the earlier and weaker justification (P1: 'AN4325's filter is recommended-only, and the mode excludes non-required filtering') with a positive engineering reason. The conclusion is unchanged - no filter - but the reason now survives a reviewer who does not accept scope arguments \\
P2 & R1 is a 10k pull-DOWN from BOOT0 to GND (boot from main flash). BRING-UP VERIFICATION ITEM: the strap polarity is not established by any allowlisted source & the P2 researcher established from the datasheet that BOOT0 has NO internal pull and so must be strapped, but correctly refused to assert WHICH level selects which memory - that mapping is in RM0360, which is on no allowlisted host, and the researcher kept it out of the record rather than smuggling it into prose. The design nonetheless needs a direction, and it rests on three converging pieces of P1 evidence: the reference-design agent read AN4325 (via a non-allowlisted mirror) showing R2 as a 10k pull-DOWN, this repo's own stm32-blinky carries 'R2 10k BOOT0 pulldown (run from flash)', and BOOT0=0 selecting main flash is uniform across the STM32 families. That is sound enough to build on and NOT sound enough to write a knowledge record about - which is exactly why it is recorded here as a decision and flagged for bring-up to confirm on hardware \\
P2 & The J2 pin-order ruling's premise - that a debug probe's rail pin is a high-impedance SENSE input, not a driver - is P1-sourced but NOT allowlist-sourced. Added to the bring-up verification list beside the BOOT0 strap polarity & the P2 swd-debug-port researcher honestly refused to assert reversal SAFETY, because it could not source the probe side from any allowlisted host (arm.com and segger.com are not on the list at all) - it grounded only the pin-order GEOMETRY and labelled the i -\textgreater{} (n+1)-i derivation as unsourced in the record prose. That is a fair challenge to my own P1 ruling and it is recorded rather than brushed aside. The evidence the ruling actually rests on came from P1 plain-web research: ST UM2448 describes the STLINK-V3SET rail pin as 'Input for STLINK-V3SET', SEGGER states J-Link's is 'not intended to power the target system', and measured VTref draw is 25-170 uA. That is sound enough to choose a pin order and not sound enough to write a knowledge record about - the same line the BOOT0 strap sits on. The board is safe either way: the ruled order is no WORSE than any alternative under any premise, and strictly better if the premise holds \\
P2 & Accepted the corrected mcu-supply-pin-pair-decoupling as ONE record at level topology bounded by vdd\_v 1.8-3.6. Did NOT take the researcher's own proposed split into a principle plus a topology record - carried to the promote pass instead & the split (unbounded principle 'the bank is counted per supply pin and per rail, not per device' + topology record holding the capacitance values with the rail bound) is genuinely the better shape and the researcher proposed it against its own interest. But shape is what the owner rules on at promotion, and the record as corrected is already sound: the uncited idd\_ma\textless{}=120 bound is gone, X5R is gone, and the rail bound is now sourced from pages the researcher personally read (AN13033 Table 6 p5 showing the bank changing DISCONTINUOUSLY below the range - 22uF+100nF+47pF on a DC-DC core rail - plus ST Table 21, with p43 added to the record's own sources). Recorded because the researcher disclosed the incentive rather than hiding it: the checklist's decoupling row is min\_level topology, so a principle-level record would have left this board showing a decoupling gap. It tested the counterfactual and said that was not its reason - that disclosure is what makes the ruling reviewable \\
P2 & The no-series-resistor ruling stands, but it is now recorded as a CONDITIONAL that this board passes, not as a family fact. Infineon AN220270 p30 recommends 20 ohm on SWCLK unconditionally ('in general, it is recommended'); this board omits it because its own tr-vs-tau check clears by \textasciitilde{}80x & the second reader flagged the tension explicitly: a future design agent reading only the record's remedy\_if\_check\_fails branch could omit a resistor Infineon would fit. The knowledge records now express it correctly - the pulls record declines to rule on series elements and hands off to the edge-rate record's per-board rise-time-against-line-delay check, and both disclose the Infineon recommendation they do not carry. So the library will not tell a future board 'no series resistor needed'; it will tell it to run the check. For bb-mcu specifically the check is decisive: 5 ns fastest specified edge (DS Table 48, and the Min column is empty on every fmax/tr/tf row so 5 ns is a specified fastest, not a guaranteed floor) against \textasciitilde{}1.5 ns settling on a \textless{}=50 mm trace and a 125 ns half-bit at 4 MHz. Flagged for the promote pass as a curation question, not a design change \\
P2 & DESIGNING UNDER COVERAGE GAP: slot block:B2 (swd-debug-port) stays gap on three classes - esd, decoupling, constraints-emission. No second research round opened, though 4 of the 6-task budget remain & these are the rows the first pass DECLARED as honest gaps rather than filling with weak sources, and a second round would hit the same wall: ESD knowledge for a debug header needs IEC 61000-4-2 or a vendor ESD app note, and the two hosts that would carry the probe-side and standards material (arm.com, segger.com) are not on the fetch allowlist at all. More important, the gap is RISKLESS on this board: esd and decoupling are feature classes the block-only scope tier EXCLUDES, so even perfect knowledge would change nothing about what gets built here - and constraints-emission at principle level has a positive answer already recorded, namely that this interface emits NO layout constraints (high\_speed and diff\_pairs carry no SWD entries, and JLC sells no impedance-controlled 2-layer stackup). The gap is a statement about the LIBRARY, not about this board. It is recorded here rather than left silent, per the build-modes rule, and is the right thing for a future MCU board at a higher scope tier to close \\
P2 & Both new checklists (mcu, swd-debug-port) and all 13 verified records sit at maturity draft/verified with floor\_met false - they are queued PENDING for the owner's promotion ruling, not promoted by this run & only owner-approved or bench-proven records satisfy coverage at the default floor, and promotion is explicitly the owner's verb (draft -\textgreater{} verified is the reader's, verified -\textgreater{} approved is the owner's, approved -\textgreater{} proven is bring-up's). This run deliberately does not self-promote: it would be the same agent grading its own homework, and the library is shared with every other board. Two queue entries are compiled in boards/bb-mcu/learnings/queue.yaml awaiting a promote pass \\
P3 & BOM keeps ST's stated capacitance values (4.7 uF bulk, 10 nF + 1 uF on VDDA) even though all three are JLC EXTENDED at 0603, rather than substituting a nearby Basic value & the part-sourcer verified with the raw ranked search window that JLC's curated Basic-passive list carries no 4.7 uF, 10 nF or 1 uF up to 1206 in this catalog snapshot (2026-08-16) - so this is a real catalog gap, not a bad query. Substituting e.g. 10 uF for the 4.7 uF bulk to reach a Basic SKU would be electrically defensible (bulk is a minimum, not a maximum) but it would depart from the value ST's Figure 12 actually states, and this board's entire discipline has been to build exactly what the datasheet requires and to make the required-vs-recommended line visible. The cost argument that would justify substituting does not even apply: the run stops at P9 by mode, so no per-unique-Extended-part assembly fee is ever incurred. 2 Basic (C1 100nF, R1 10k) + 6 Extended; \textasciitilde{}1.35 USD/board, \textasciitilde{}6.76 USD for the 5-board build \\
P3 & U1 alternate is C32908 (STM32F030F4P6, tube pack, stock 1329) - same die and pinout, differing only in packaging from C89040's tape-and-reel & the P1 scout flagged every candidate as single-source at its exact LCSC row. This is the closest thing to a true drop-in the catalog offers (LCSC serves the identical datasheet document for both). The PY32F030 is deliberately NOT listed as an alternate despite being pin-compatible: it was rejected at P1 on the hard ground that puya.com is not on the research fetch allowlist, so its behaviour can never be sourced by this pipeline, and it is not a verified drop-in. U1 stock re-verified live at 18,977 (19,015 at P1 - normal drift) \\
P3 & ACCEPTED the pulled U1 SOP-20 land pattern despite two deviations from ST Figure 37: pad 0.35 x 1.494 mm vs 0.40 x 1.35 mm, and row spacing 6.00 mm vs 5.75 mm. No hand edit to the footprint & worked the geometry rather than judging the delta. At the package's 6.4 mm lead span with a \textasciitilde{}0.6 mm foot, the lead foot sits about 2.6-3.2 mm from the centreline. ST's land spans 2.20-3.55 mm from centreline; the pulled land spans 2.253-3.747 mm. BOTH fully capture the foot - the pulled one trades \textasciitilde{}0.05 mm of heel for \textasciitilde{}0.20 mm more toe, which is the direction that helps fillet formation and inspection. The narrower pad also WIDENS the inter-pad gap from 0.25 mm (ST) to 0.30 mm, so solder-bridging risk goes DOWN, not up. It is the EasyEDA/LCSC land for this exact LCSC part number, i.e. the one JLC's own PCBA line places. Verified pin-1 independently: silk dot and chamfered corner both anchor at pad 1 only, confirmed in a kicad-cli render \\
P3 & DEFECT FOUND AND FIXED: J1's pulled footprint (C8465, WJ500V-5.08-2P) drills 1.30 mm where the vendor's own PCB LAYOUT panel recommends 1.50 mm. Approving a hand edit to drill 1.50 mm with the pad enlarged to 2.30 mm (annulus 0.40 mm) & the datasheet dimensions the pin at 0.90 mm wide. If it is square - which the square screw-clamp opening drawn above it indicates, though the shape is not explicitly labelled - the worst-case diagonal is 0.90{*}sqrt(2) = 1.273 mm, which leaves 0.027 mm of total diametral clearance in a 1.30 mm hole. That is a terminal that does not seat, found at assembly on all 5 boards. The vendor's recommended 1.50 mm is authoritative whichever shape the pin is, so the fix does not depend on resolving the ambiguity. Enlarging the pad from 2.00 to 2.30 mm keeps the annulus at 0.40 mm instead of dropping it to 0.25 mm: a screw terminal takes real mechanical stress from wire pull, so barrel strength is worth the copper, and at 5.08 mm pitch a 2.30 mm pad still leaves 2.78 mm of clearance. This is the exact class of defect the librarian role calls a top-3 board killer, and NOTHING in the pipeline would have caught it - fp\_verify had no land pattern to diff against until this extraction was made, and it does not check drill-vs-pin at all. The J2/J3 header needs no change: 1.10 mm pulled vs 1.02 mm recommended is oversize, the safe direction, with 0.195 mm clearance on its SQ 0.64 mm pin \\
P3 & P3-exit coverage: all 3 part slots COVERED (C89040, C8465, C32713271); block:B1 provisional; block:B2 still gap on the same esd/decoupling/constraints-emission already recorded at P2. NO new research tasks opened at P3 & the escalated first call (--maturity-floor proven --research-provisional, mandatory at a learning target) reported all five slots as gaps, but that is the escalation re-firing on slots this run ALREADY researched at P2, not a new finding: with the floor at proven, a verified-but-unapproved record can never satisfy anything until a built board proves it via knowledge.py --prove, so a seeding run will always show block gaps on that call. The re-run without the flags gives the true verdict and it is identical to P2's post-research state. Opening tasks would re-research what was researched two hours ago. Side effect worth noting: the two connector datasheet extractions ordered to chase the J1 drill defect also closed the part-slot coverage - C8465 and C32713271 went from having no extraction at all to covered \\
P4 & ONE schematic-block agent for the whole board, not one per sheet & the P2 sheet plan is a single flat root sheet for 9 parts and 10 nets, and the full-run recipe explicitly sanctions one schematic agent on a small board provided it is recorded. Splitting it would add a stitching seam with nothing on either side of it, and the recorded /sheet/LABEL export hazard that cost lumina-par is avoided entirely by having no hierarchy \\
P4 & REVERSING the P1 ruling: ADD C5, 100 nF from NRST to GND. The reviewer was right and my original reason was wrong & I ruled it out at P1 on AN4325 filing the part under Table 6 Optional components 'for RESET button', reasoning that this board has no button. The reviewer showed that misreads the evidence: the DATASHEET's own Figure 21 note 1 recommends the capacitor unconditionally as parasitic-reset protection, which is a different and stronger justification than button debounce, and the run's own C89040 extraction and its own swd-debug-port research record both carry the 0.1 uF. Decisive for scope: requirements.md section 2 point 5, written at P0 before any of this, explicitly scoped 'whatever the datasheet requires on NRST / boot-mode pins' as in-scope support - so this is not the recommended-only conditioning the block-only tier excludes, it is the carve-out the requirements anticipated. And the exposure is specific to this board: NRST is deliberately routed to an unshielded flying lead, which is exactly the parasitic-reset case the note addresses. My own P1 decision text had already flagged that risk and accepted it; the reviewer supplied the source that turns the accepted risk into an unnecessary one \\
P4 & J3 pin order CHANGED from IO1/IO2/GND/IO3/IO4 to IO1/IO2/IO3/IO4/GND - GND moves off the centre to position 5 & the reviewer found that J2 and J3 are the SAME unkeyed 5-position part, so a cable made for one fits the other. With GND at J3 centre and +3V3 at J2 centre, plugging the GPIO cable into the debug header shorts the rail through the cable's own return - a hard short, the worst outcome available on this board. Moving GND to an end makes the worst cross-plug outcome a GPIO pin contacting the rail or ground, both current-limited by the pin's own driver and non-destructive at 3.3 V. The cost is the return-path benefit the P2 architect chose centre-GND for, which is immaterial here: the owner's answer 3 freed these four pins from any electrical requirement, they are plain digital I/O with no stated edge rate, and a scope probe brings its own ground clip. Accepted trade: J3's own 180-degree reversal is now less clean (GND no longer self-maps), but reversal costs an IO-to-GND contact while cross-plugging cost a rail short \\
P4 & J2 pin order UNCHANGED (GND/SWCLK/3V3/SWDIO/NRST) despite the reviewer's one-position-slip warning. Mitigation stays per-pin silk, carried to P6 & the reviewer correctly notes my order was ruled on the 180-degree reversal case only, and that a one-position slip puts a probe push-pull driver onto the rail. I worked the alternatives and centre-rail still wins: with the rail at an END, a one-position slip lands the probe's GND wire on it, which is a hard short - strictly worse than a driver contact, which is current-limited by the driver and which probes generally detect. There is no 5-position order that is safe against BOTH a reversal and a one-position slip, so the ruling is to take the failure that is current-limited in every case. Recorded rather than dismissed: the residual risk is real and its only mitigation is that the per-pin silk actually gets placed and stays legible, which constraints.json already flags as load-bearing and which P6 must deliver \\
P6 & PIPELINE DEFECT: at a geometry-OUTPUT binding the place gate must run AFTER board\_edit --outline fit, not before. The full-run recipe documents the opposite order and it cannot pass & placelib.legality\_violations measures every declared placement.edges ref against the CURRENT outline with EDGE\_TOL 2.5 mm. At a canonical binding the provisional outline is deliberately oversized, so a canonically placed connector sits far from it - here J1 was 16.9 mm from the provisional right edge and the gate failed on it. The only way to satisfy the edge check against a provisional outline is to spread the connectors out to it, after which --outline fit measures the spread and 'earns' the provisional size. That is exactly the bb-buck failure with the sign flipped, and the U18 mechanism defeats itself. The placement agent placed canonically anyway, refused to resize, and machine-proved the fix is the ORDER: fit first at any margin at or under 2.5 mm, then gate, which passes 0/0. Fitting first is also better on the merits - the gate then judges the geometry that will actually be built. Recorded as a defect in the recipe, not in this board \\
P6 & Board size EARNED: 34.77 x 22.26 mm (774 mm2), down from the provisional 51.15 x 43.52 (2226 mm2). Hand-built canonical layout accepted over both annealer candidates & geometry is an OUTPUT at this binding and this is the output. The annealer returned one legal candidate at 200.12 mm HPWL with only 3 movable clusters, so the SA numbers are noise on a board this small; the hand layout reaches 161.36 mm with the same zero legality violations, and route\_auto --probe completes 1.00 with 0 unrouted on the fitted geometry. U1 at rot 270 is the whole floorplan: the SWD end faces J2/top, the VDD/VSS face faces J1/right, and the PA0-PA3 face faces J3/left, so all three escapes are non-crossing. The seed had U1 at rot 0, which crossed the GPIO escapes under the package \\
P6 & ACCEPTED the placement agent's overrule of 'never move board\_only refs': it moved H1-H4, and rejected four TRUE corner holes in favour of a top pair at the corners with the bottom pair inset & the overrule was correct and correctly flagged. board\_edit --outline fit includes mounting-hole courtyards in the content bbox, so leaving H1-H4 at the provisional corners would have frozen the board at 51 x 43.5 mm - the holes would have defeated the fit. constraints.json also assigns the hole count and placement to P6 explicitly, and requirements s5 sanctions fewer than four if the earned outline cannot hold them. On the count: four true corner holes force about 33.8 x 26.5 mm, roughly +25 percent area and \textasciitilde{}9 mm of pure air, to satisfy a board with no enclosure and no mechanical requirement beyond mounting the bare board - that is the padding the block-only tier exists to refuse. Residual accepted: the rightmost \textasciitilde{}10 mm of the bottom edge is unsupported, but J1 is braced by H2 above and H3 below-left, so the screw terminal is not cantilevered off an unsupported corner \\
P6 & J1 rotation 270 PROVEN by orthographic side render, and the finding contradicts this repo's existing connector learning: WJ500V-5.08-2P's wire mouth is at footprint-local -Y, the REVERSE of the KF128/DB128L family & the role prompt and the constraints note both refuse to guess a footprint's native entry direction because a prior board recorded the cost of guessing it. The agent proved this one instead of assuming, and found the top view reads BACKWARDS on this part. Had it carried the KF128 assumption across, the screw terminal's wire opening would have faced INTO the board and the board would have been unusable at the bench - a defect no gate checks, because entry direction is not in any machine-readable field. Silk verified on the fitted render by me directly: +3V3 above the top pole, -GND below the bottom one, each adjacent to its own pole and 9 mm from the other \\
P6 & J2 and J3 per-pin silk VERIFIED on the fitted top render, closing the P4 review's hard requirement & the P4 review established that per-pin silk is the ONLY mitigation for the two unkeyed-header hazards (180 degree reversal, and cross-plugging the GPIO cable into the debug header). Checked on the render rather than taken on report: J2 carries a pin-1 marker at the right and reads right-to-left GND / SWCLK / 3V3 / SWDIO / NRST, matching the ruled order, using a two-row stagger because 5-character names do not fit a single row at 2.54 mm pitch. J3 carries one label per pin plus a 1 at the top and reads IO1 / IO2 / IO3 / IO4 / GND, also the ruled order. Both legible at render scale \\
P8 & FORCED TIER DOWNGRADE: the P8 verify-reviewer runs on sonnet, below its fable/high spec and below the opus/xhigh substitution used all run. Recorded, not silent & the account hit its weekly usage limit mid-review (opus exhausted; fable had already been unavailable since P0). The playbook's rule is to substitute one effort step UP and never silently drop to a weaker tier - I cannot substitute up, so I am recording the drop instead of hiding it. Mitigation: the machine state this reviewer backstops is fully green with no skips (erc 0/0, place 0/0, drc\_routed 0/0, verify 0/0 across all 8 checks, irdrop and pdn\_z pass), so the reviewer is hunting only what scripts cannot see on a 10-part 2-layer board - the least demanding review surface this pipeline produces. The first partial pass, before it died, had already independently confirmed the highest-stakes item: J1's wire mouth faces OFF-board, verified in a +X orthographic view where both funnels land exactly on the two pole positions. If the owner wants the review re-run at full tier after the limit resets, the board is committed at da26db0 and re-running is a single spawn \\
P8 & P8 review F1 (no mounting hole near J1's bottom-right corner) ACCEPTED as a warning, no change. Consistent with the P6 ruling & the reviewer's number is sharper than mine at P6 - H3 is inset 14.3 mm from the bottom-right corner against H1/H2's 4 mm true-corner inset - but the conclusion holds. J1 at (37.8, 21.9) is braced diagonally on two sides: H2 at (39.85, 12.94) is 9.2 mm above-right and H3 at (29.55, 26.1) is 9.2 mm below-left, so the terminal is not cantilevered off an unsupported corner. The board is 1.6 mm FR4 at 34.77 x 22.26 mm, which is stiff at that span. Adding a fifth hole or moving H3 outward costs board area on a tier that excludes mechanical features beyond mounting the bare board, to harden a bench article against a load nobody has specified. The reviewer itself files this as a bench-reliability question rather than a defect, and its own OPEN says the severity depends on handling the run cannot know \\
P8 & P8 review F2 (0.582 mm2 zero-anchor copper island on B.Cu) carried to P9 DFM undecided, deliberately. Not fixed at P8 & the reviewer independently re-ran plane\_repair --flag-only rather than taking my report, confirmed 0 anchors, and recommended DFM purge it since zero anchors means zero risk in removing it. But its own OPEN is the reason not to act yet: whether a 0.34 mm gap survives JLC's process is a fab-process question that needs the actual DFM pass, not a render. Fixing it now would re-stale drc\_routed and verify for a change DFM may render moot or may specify differently. If the DFM leg flags it, it becomes a work order with a known-safe fix; if DFM is clean, the island is electrically inert on a 3.3 V bench board and stays. Deciding it at P9 with evidence beats deciding it at P8 without \\
P9 & P8 review F2 RESOLVED by evidence: the 0.582 mm2 zero-anchor copper island STAYS. The DFM gate did not flag it, with all 8 legs running & this is the decision I deliberately deferred from P8 rather than making it blind. The dfm gate passed 0/0 with copper, copper\_to\_edge, drill, hole\_to\_edge, silk, polarity, bom and release all RAN and passed - no skips, no waivers. The island is 0.34 mm from the main pour and carries zero anchors, so it is electrically inert on a 3.3 V bench board, and the fab-process question the reviewer correctly said needed a real DFM pass has now been answered by one. Removing it would have re-staled drc\_routed and verify for no measured benefit \\
P9 & parts.json is now placed BESIDE the board at kicad/parts.json, alongside constraints.json and decoupling.json & acting on a prior run's recorded learning: gate.py --gate dfm resolves parts.json from the BOARD's directory, so the BOM-completeness leg silently never runs when the file lives only at parts/parts.json - a pass that means nothing. Verified the fix worked rather than assuming: the gate's coverage block shows the bom leg in ran and passed, so it genuinely executed. Sidecars belong beside the board from P5 onward and parts.json is one of them \\
\bottomrule
\end{longtable}
}
\subsection*{Sheet plan}
\subsection*{sheets.md - bb-mcu sheet plan, refdes ranges, CANONICAL net names}
\subsubsection*{1. Sheet plan: ONE FLAT ROOT SHEET}
\begin{quote}\small\ttfamily\raggedright
\textbar{} Sheet \textbar{} File \textbar{} Blocks it carries \textbar{} Interface nets at its boundary \textbar{} Refdes ranges \textbar{} `pwr\_base` \textbar{}\\
\textbar{}---\textbar{}---\textbar{}---\textbar{}---\textbar{}---\textbar{}---\textbar{}\\
\textbar{} root (only) \textbar{} `kicad/bb-mcu.kicad\_sch` \textbar{} B1 MCU minimum system, B2 SWD debug port, J1 power in, J3 GPIO out, H1-H4 mounting \textbar{} none - the board's only interfaces are the three physical connectors \textbar{} `U1-U9`, `C1-C19`, `R1-R9`, `J1-J9`, `H1-H9` \textbar{} {*}{*}1{*}{*} (`\#PWR01..`, `\#FLG01..`) \textbar{}
\end{quote}
\textbf{A hierarchy is REJECTED, explicitly and with the argument stated}, because the role prompt invites disagreement and there is none to have here.

This board has \textbf{nine electrical parts and ten nets} on one A4 page. Hierarchy exists to keep large schematics readable and to give each sheet its own refdes and \texttt{\#PWR} range; neither problem exists at this size. The cost of splitting is concrete and this repo has already paid it: a label inside a child sheet exports as \texttt{/\textless{}sheet\textgreater{}/\textless{}LABEL\textgreater{}}, which silently unhooks every \texttt{constraints.json} entry that spells it \texttt{/\textless{}LABEL\textgreater{}} - the P4 amendment class that cost \texttt{lumina-par}, and the reason \texttt{sbuck-5v3a} and \texttt{bb-buck} both went flat. With one root sheet there are no sheet prefixes anywhere, by construction.

Refdes ranges are therefore trivially unique. They are stated anyway because the rule is "unique ACROSS sheets": {*}{*}if a later amendment ever splits this sheet, the child sheets take \texttt{pwr\_base = 40} and \texttt{80} and keep the prefix ranges above{*}{*}, and every net name in s2 must be re-checked against the new export.

Mounting holes are reserved \texttt{H1-H4}. {*}{*}Whether four or two are fitted is a P6 call{*}{*}, not a P2 one: \texttt{requirements.md} s5 asks for four M3 clearance holes inset from the corners and sanctions two on opposite corners if the earned outline cannot hold four. Under \texttt{canonical} binding the outline is not known here, so the count is settled when it is known and the choice is recorded.

\subsubsection*{2. CANONICAL NET NAMES - binding from here}
\texttt{research/interface-swd.json} proposed \texttt{SWDIO / SWCLK / NRST / +3V3 / GND} as "proposals in the standard's conventional form" and handed reconciliation to P2. \textbf{These are now the names the schematic MUST produce}, and they are the names \texttt{constraints.json} uses.

\begin{quote}\small\ttfamily\raggedright
\textbar{} Canonical net \textbar{} KiCad mechanism \textbar{} Spans \textbar{} In `constraints.json` \textbar{}\\
\textbar{}---\textbar{}---\textbar{}---\textbar{}---\textbar{}\\
\textbar{} {*}{*}`+3V3`{*}{*} \textbar{} power symbol (VALUE = net name) -\textgreater{} exports BARE \textbar{} J1.1 -\textgreater{} C1 C2 C3 C4 -\textgreater{} U1 VDD(16) + VDDA(5) -\textgreater{} J2.3 \textbar{} `power` 0.1 A \textbar{}\\
\textbar{} {*}{*}`GND`{*}{*} \textbar{} power symbol -\textgreater{} exports BARE \textbar{} everything; B.Cu pour \textbar{} {*}{*}not declared{*}{*} - see `power\_tree.md` s7 \textbar{}\\
\textbar{} {*}{*}`/SWDIO`{*}{*} \textbar{} root-sheet LOCAL LABEL -\textgreater{} exports with ONE leading slash \textbar{} U1 pin 19 (PA13) -\textgreater{} J2.4 \textbar{} not declared (advisory only, `notes`) \textbar{}\\
\textbar{} {*}{*}`/SWCLK`{*}{*} \textbar{} root-sheet local label \textbar{} U1 pin 20 (PA14) -\textgreater{} J2.2 \textbar{} not declared (advisory only, `notes`) \textbar{}\\
\textbar{} {*}{*}`/NRST`{*}{*} \textbar{} root-sheet local label \textbar{} U1 pin 4 -\textgreater{} J2.5 \textbar{} not declared \textbar{}\\
\textbar{} {*}{*}`/BOOT0`{*}{*} \textbar{} root-sheet local label \textbar{} U1 pin 1 -\textgreater{} R1 \textbar{} not declared \textbar{}\\
\textbar{} {*}{*}`/IO1`{*}{*} \textbar{} root-sheet local label \textbar{} U1 pin 6 (PA0) -\textgreater{} J3.1 \textbar{} not declared \textbar{}\\
\textbar{} {*}{*}`/IO2`{*}{*} \textbar{} root-sheet local label \textbar{} U1 pin 7 (PA1) -\textgreater{} J3.2 \textbar{} not declared \textbar{}\\
\textbar{} {*}{*}`/IO3`{*}{*} \textbar{} root-sheet local label \textbar{} U1 pin 8 (PA2) -\textgreater{} J3.4 \textbar{} not declared \textbar{}\\
\textbar{} {*}{*}`/IO4`{*}{*} \textbar{} root-sheet local label \textbar{} U1 pin 9 (PA3) -\textgreater{} J3.5 \textbar{} not declared \textbar{}
\end{quote}
Why \texttt{+3V3} and not \texttt{3V3}: it is a rail feeding supply pins and decoupled by C1-C4, so it is drawn with a power symbol, and a power symbol's exported name is bare. Why \texttt{/IO1} and not \texttt{IO1}: root-sheet local labels are prefixed with one slash on export. \textbf{A power symbol WINS over a coincident label}, so P4 must not mix the two mechanisms on one node.

\textbf{Silk drops the slash.} The header silk reads \texttt{IO1 IO2 GND IO3 IO4} and \texttt{GND SWCLK 3V3 SWDIO NRST} - the slash is a netlist artifact, not a label a bench user should ever see. Naming note from the research: ST never calls the sense pin "VTREF" on its own connectors (it is \texttt{T\_VCC} or \texttt{VDD\_TARGET}), so the silk says \texttt{3V3}.

\textbf{Enforcement, not hope:} at P4 run \texttt{netlist\_audit.py --net kicad/bb-mcu.net --constraints architecture/constraints.json}. A net spelled differently in the schematic surfaces as \texttt{missing\_net} (error). Without that run a misspelling is silent - \texttt{check\_current} finds no copper on a net that does not exist and reports nothing.

\subsubsection*{3. Contents of the root sheet}
\paragraph*{U1 - STM32F030F4P6TR, TSSOP-20 - complete pin map}
\begin{quote}\small\ttfamily\raggedright
\textbar{} Pin \textbar{} Name \textbar{} Net \textbar{} Note \textbar{}\\
\textbar{}---\textbar{}---\textbar{}---\textbar{}---\textbar{}\\
\textbar{} 1 \textbar{} BOOT0 \textbar{} `/BOOT0` \textbar{} -\textgreater{} R1 -\textgreater{} `GND`. Dedicated pin, NO internal pull, hardware-sampled on the 4th SYSCLK edge after reset \textbar{}\\
\textbar{} 2 \textbar{} PF0-OSC\_IN \textbar{} - \textbar{} unused, left floating (no crystal; tie-offs are RECOMMENDED only and mode-excluded) \textbar{}\\
\textbar{} 3 \textbar{} PF1-OSC\_OUT \textbar{} - \textbar{} unused, left floating \textbar{}\\
\textbar{} 4 \textbar{} NRST \textbar{} `/NRST` \textbar{} -\textgreater{} J2.5. Permanent internal pull-up 25/40/55 k; NOTHING added \textbar{}\\
\textbar{} 5 \textbar{} VDDA \textbar{} `+3V3` \textbar{} tied to VDD by a BARE TRACE; C3 10 nF + C4 1 uF to `GND` \textbar{}\\
\textbar{} 6 \textbar{} PA0 \textbar{} `/IO1` \textbar{} -\textgreater{} J3.1 \textbar{}\\
\textbar{} 7 \textbar{} PA1 \textbar{} `/IO2` \textbar{} -\textgreater{} J3.2 \textbar{}\\
\textbar{} 8 \textbar{} PA2 \textbar{} `/IO3` \textbar{} -\textgreater{} J3.4 \textbar{}\\
\textbar{} 9 \textbar{} PA3 \textbar{} `/IO4` \textbar{} -\textgreater{} J3.5 \textbar{}\\
\textbar{} 10 \textbar{} PA4 \textbar{} - \textbar{} unused \textbar{}\\
\textbar{} 11 \textbar{} PA5 \textbar{} - \textbar{} unused \textbar{}\\
\textbar{} 12 \textbar{} PA6 \textbar{} - \textbar{} unused \textbar{}\\
\textbar{} 13 \textbar{} PA7 \textbar{} - \textbar{} unused \textbar{}\\
\textbar{} 14 \textbar{} PB1 \textbar{} - \textbar{} unused \textbar{}\\
\textbar{} 15 \textbar{} VSS \textbar{} `GND` \textbar{} the only ground pin; also the VDDA return (no VSSA on this package) \textbar{}\\
\textbar{} 16 \textbar{} VDD \textbar{} `+3V3` \textbar{} C1 100 nF + C2 4.7 uF \textbar{}\\
\textbar{} 17 \textbar{} PA9 \textbar{} - \textbar{} unused \textbar{}\\
\textbar{} 18 \textbar{} PA10 \textbar{} - \textbar{} unused \textbar{}\\
\textbar{} 19 \textbar{} PA13 \textbar{} `/SWDIO` \textbar{} -\textgreater{} J2.4. Dedicated SWD after reset, internal PULL-UP active \textbar{}\\
\textbar{} 20 \textbar{} PA14 \textbar{} `/SWCLK` \textbar{} -\textgreater{} J2.2. Dedicated SWD after reset, internal PULL-DOWN active \textbar{}
\end{quote}
Nine pins are deliberately unconnected. {*}{*}PA13/PA14 must NOT also be routed to J3{*}{*} - they can be released to GPIO in software, but J2 needs them as SWD permanently and routing them to both headers is a pin conflict.

\paragraph*{The rest of the board}
\begin{quote}\small\ttfamily\raggedright
\textbar{} Refdes \textbar{} Part \textbar{} Net attachments \textbar{}\\
\textbar{}---\textbar{}---\textbar{}---\textbar{}\\
\textbar{} `C1` \textbar{} 100 nF X7R ceramic \textbar{} `+3V3` / `GND`, at U1 pins 16/15 \textbar{}\\
\textbar{} `C2` \textbar{} 4.7 uF X7R ceramic \textbar{} `+3V3` / `GND`, bulk on the VDD net \textbar{}\\
\textbar{} `C3` \textbar{} 10 nF X7R ceramic \textbar{} `+3V3` / `GND`, at U1 pin 5 \textbar{}\\
\textbar{} `C4` \textbar{} 1 uF X7R ceramic \textbar{} `+3V3` / `GND`, at U1 pin 5 \textbar{}\\
\textbar{} `R1` \textbar{} 10 k (5 \% is fine) \textbar{} `/BOOT0` -\textgreater{} `GND` \textbar{}\\
\textbar{} `J1` \textbar{} 2-pole 5.08 mm screw terminal, THT \textbar{} 1 = `+3V3`, 2 = `GND`. Silk `+` / `-` \textbar{}\\
\textbar{} `J2` \textbar{} 1x5 0.1 in header, single row, straight, THT \textbar{} 1 `GND`, 2 `/SWCLK`, 3 `+3V3`, 4 `/SWDIO`, 5 `/NRST` \textbar{}\\
\textbar{} `J3` \textbar{} 1x5 0.1 in header, single row, straight, THT \textbar{} 1 `/IO1`, 2 `/IO2`, 3 `GND`, 4 `/IO3`, 5 `/IO4` \textbar{}\\
\textbar{} `H1`-`H4` \textbar{} M3 clearance holes, 3.2 mm \textbar{} mechanical, no net (count settled at P6) \textbar{}
\end{quote}
\textbf{Voltage ratings}: everything sits on a 3.3 V rail, so 16 V or 25 V X7R ceramics are ample and DC-bias derating is a non-issue at 3.3 V on parts rated 5-8x that. P3 picks the exact parts; nothing here constrains them beyond value, dielectric and a package the JLC PCBA line places on the top side.

\paragraph*{J2 pin order - RULED at P1, do not re-order}
\texttt{GND / SWCLK / 3V3 / SWDIO / NRST}, positions 1..5. The reasoning is in \texttt{blocks.md} s2 and in \texttt{state.json}; the one-line version is that on an unkeyed 5-way shell a 180-degree reversal maps position \texttt{i} to \texttt{6-i}, so {*}{*}3V3 at the centre is the unique arrangement in which a reversed probe lands its high-Z VTref INPUT on the board's rail rather than a probe OUTPUT.{*}{*} This is deliberately NOT a vendor cable order - flying-lead probes present individually labelled wires, so per-pin silk is what the user reads either way.

\paragraph*{J3 pin order - RULED HERE at P2}
\texttt{IO1 / IO2 / GND / IO3 / IO4}, positions 1..5. The owner fixed the SET (four GPIO plus a ground, because a signal with no return is not usable at the bench); the ORDER was open. Ground goes in the CENTRE, for three reasons:

1. \textbf{Return distance.} With GND at an end, \texttt{IO4} is four positions from its return; at the centre no signal is more than two positions away. On a bench header the ground lead IS the return path, and SEGGER's own note is that "long ground leads can significantly worsen ringing and distortion". This is the real reason. 2. \textbf{The same reversal mechanism that decided J2.} Five positions, unkeyed, \texttt{i -\textgreater{} 6-i}; a centre GND maps to itself, so a reversed plug still lands ground on ground. Nothing destructive can happen on J3 either way (there is no rail on this header), but a user clipping a ground lead and getting a driven GPIO is a real bench annoyance that costs nothing to prevent. 3. \textbf{Consistency}: both headers on this board answer "where does ground go?" the same way, which is the kind of thing a study article should teach once rather than twice differently.

\subsubsection*{4. Constraint keys deliberately ABSENT - read before "fixing" them}
\begin{quote}\small\ttfamily\raggedright
\textbar{} Key \textbar{} Status \textbar{} Why \textbar{}\\
\textbar{}---\textbar{}---\textbar{}---\textbar{}\\
\textbar{} `diff\_pairs` \textbar{} {*}{*}OMITTED entirely{*}{*} - NOT `{[}{]}` \textbar{} P1 ruling. An explicit empty list DISABLES `check\_diffpair` board-wide; omitting the key leaves auto-discovery ARMED. Both are no-ops on a board with no pair, so take the fail-safe one - disabling a check board-wide to say "we have none" is how a silent blind spot gets built in for the next edit. None of this board's net names matches the auto-discovery suffixes (`\_P`/`\_N`, `DP`/`DM`, `D+`/`D-`), so nothing is discovered by accident. This DIVERGES from `research/interface-swd.json`, which emitted `{[}{]}`, and from `stm32-blinky`, which shipped `{[}{]}`. \textbar{}\\
\textbar{} `high\_speed` \textbar{} absent \textbar{} SWD binds nothing (`blocks.md` s3). Declaring these nets on 2 layers would only manufacture `corridor\_void` findings that are unfixable by construction. The `\textless{}= 50 mm` / over-the-pour / keep-SWCLK-clear guidance lives in `notes` as ADVISORY. \textbar{}\\
\textbar{} `voltages` \textbar{} absent \textbar{} Nothing on the board exceeds 3.3 V. `check\_creepage`'s derived check engages only ABOVE 30 V, so entries at 3.3 V would be dead weight, not forward-compatibility. `requirements.md` s8: the \textgreater{}30 V flag DOES NOT APPLY. \textbar{}\\
\textbar{} `voltage\_pairs` \textbar{} absent \textbar{} No bridge, no AC tap, no net pair with a differential the node voltages cannot express. \textbar{}\\
\textbar{} `thermal` \textbar{} absent \textbar{} 0.33 W at the budget ceiling, no exposed pad, no via array (`power\_tree.md` s4). \textbar{}\\
\textbar{} `planes` \textbar{} absent \textbar{} The 2-layer default (B.Cu GND pour) is exactly right, and a `planes{[}{]}` list REPLACES the defaults entirely (`stackup.md` s4). \textbar{}\\
\textbar{} `power` entry for `GND` \textbar{} absent \textbar{} Would manufacture a `check\_pdn` `pdn\_undecoupled` error requiring a waiver, and buy nothing at 0.1 A (`power\_tree.md` s7). \textbar{}\\
\textbar{} `placement.keepouts` \textbar{} absent \textbar{} Keepouts are BOARD-LOCAL coordinates and this board has no outline yet by design. The mounting-hole washer keepouts get their rects at P6, once the outline is earned. \textbar{}\\
\textbar{} `placement.separation` \textbar{} absent \textbar{} Centre-to-centre, and SKIPPED when either ref is locked - it cannot express anything this board needs. Nothing here needs parts held apart. \textbar{}\\
\textbar{} `placement.fixed` \textbar{} absent \textbar{} Nothing has an earned position at P2. Locking refs is a P6 act, after explicit `place\_edit`. \textbar{}
\end{quote}
\subsubsection*{5. Placement groups this sheet implies}
One group, \texttt{mcu}: anchor \texttt{U1}, members \texttt{C1 C2 C3 C4 R1}.

\begin{itemize}
\item \textbf{C1 (100 nF) closest to pins 16/15}, then C2 (4.7 uF) as the reservoir
behind it; \textbf{C3 (10 nF) closest to pin 5}, then C4 (1 uF). Each cap's ground
drops straight to the B.Cu pour through its own via, with the supply via and
the ground via straddling the cap right at the pin (AN4325 Fig 8). The
datasheet's word is "must": "as close as possible to, or below, the
appropriate pins".
\item \textbf{R1 in the pin-1 window} - the BOOT0 node is hardware-sampled before any
software runs and there is no reason for it to be long.
\item \textbf{The one congestion point}: C3/C4 sit at VDDA (pin 5) on the same package
face the four GPIO escape past. Place them toward the pin-1 end of pin 5's
escape so the far half of that face stays clear (\texttt{blocks.md} s4).
\end{itemize}
Connector edges are in \texttt{constraints.json} \texttt{placement.edges}: J1 on its own edge, J2 and J3 on two others, which is \texttt{requirements.md} s5's requirement that power not share an edge with the two signal headers. {*}{*}Every connector's opening or mating face must point OFF-BOARD{*}{*} - a requirement in words, because no schema key expresses it, and \texttt{rot} is deliberately left unset since the footprints' native entry directions do not exist yet (\texttt{bb-buck} recorded the cost of guessing them).

Orientation guidance for P6, not a constraint: put U1's {*}{*}pin-1 end toward J2's edge{*}{*}, so \texttt{/SWDIO}, \texttt{/SWCLK} and \texttt{/NRST} all escape from the same package end, and the four GPIO escape from the diagonally opposite region toward J3.

Full architecture narrative (not inlined; see the artifact index): \texttt{architecture/blocks.md}, \texttt{architecture/power\_tree.md}, \texttt{architecture/stackup.md}.
\section{Schematic}
Gate \texttt{erc} (phase P4): \textbf{pass} after 2 attempt(s); last 2026-08-16T21:57:21, failing 0 of 0.
The full schematic PDF follows.
\includepdf[pages=-]{reports/schematic.pdf}
\section{Layout}
Gate \texttt{place} (phase P6): \textbf{pass} after 2 attempt(s); last 2026-08-16T22:57:35, failing 0 of 0.

Gate \texttt{drc\_routed} (phase P7): \textbf{pass} after 3 attempt(s); last 2026-08-16T23:20:27, failing 0 of 0.
\emph{No board renders found.}
\section{Verification}
Gate \texttt{verify} (phase P8): \textbf{pass} after 2 attempt(s); last 2026-08-16T23:20:27, failing 0 of 0.
\emph{reports/verify\_all.json not available.}
\subsection*{Design review of record}
\subsection*{Board review - bb-mcu (adversarial, P8)}
Reviewer: verify-reviewer subagent, fresh context. All eight machine checks plus \texttt{drc\_routed}, \texttt{place}, \texttt{erc} and \texttt{verify} are clean (0/0 throughout, no skips) - taken as given, not re-litigated. This pass hunts what those scripts cannot see: renders I made myself (`reports/review-renders/bb-mcu\_\{top,bottom,iso,left, right,front,back\}.png`, 2400 px), direct extraction of pad/net/track/zone geometry from the built \texttt{.kicad\_pcb} via the KiCad 10 Python API, and an independent re-run of \texttt{plane\_repair.py --flag-only} against the built board.

\textbf{Scope tier applied}: \texttt{learning block-basics:} -\textgreater{} \texttt{block-only}, binding \texttt{canonical}. Per \texttt{reference/build-modes.md} this excludes protection, filtering, indicators, test-points, config, second-rail, mechanical and enclosure-fit - none of that is reported below. Geometry is an OUTPUT the board earned (34.77 x 22.26 mm from a 51.15 x 43.52 provisional outline, recorded P6 decision); the board is not compared to that provisional number.

\textbf{Gate: 0 errors / 2 warnings.}

---

\subsubsection*{Findings, worst first}
\paragraph*{F1 (warning) - the bottom-right corner, where the wire-tug loads land, has no mounting hole nearby}
\texttt{kind: unsupported-connector-corner} \textbar{} domain: \texttt{placement} \textbar{} refs J1, H2, H3

Measured from the built board (bbox 8.98,8.939 to 43.85,31.3 mm): H1 and H2 (top pair) sit at a true 4.0 mm inset from BOTH nearest edges - genuine corners. H3 and H4 (bottom pair) do not: H4 is 8.2 mm in from the left edge, and \textbf{H3 is 14.3 mm in from the right edge}, because J3 and J1 occupy the bottom-left and right-edge real estate respectively (\texttt{state.json} P6 decision, correctly earning the board this shape). That leaves the entire bottom-right corner region - exactly where the screw terminal J1 sits - with no anchor point closer than 9.2 mm (J1 to H2 measures 9.19 mm, J1 to H3 measures 9.26 mm; both are ordinary distances to the NEAREST hole, but neither hole pins the corner itself, so that corner can flex/lift under a lever load in a way a true corner hole would resist).

J1 is not a passive connector here - it is the interface that gets a wire pushed in, a wire pulled out, and two screws torqued down, repeatedly, at the bench, every time this board is used (\texttt{requirements.md} s1: "let it be programmed and debugged" implies re-plugging cycles; the terminal itself is rated for wire pull per its datasheet, but the PCB mounting is what resists that force being transmitted into the board). Concentrating the least-anchored corner directly under the highest cyclic mechanical load is the kind of thing DRC/ERC/placement gates cannot see, because "four M3 holes present" is satisfied literally (\texttt{gate-place} passed `courtyard/outline/edges/keepouts/ decoupler\_distance`, none of which score corner support quality).

Not a fabrication defect and not a call to change the earned geometry - FR4 is stiff at this scale and 9.2 mm is not a dramatic cantilever, so this is not filed as an error. But it is a real, board-specific bench-reliability question or the human, worth weighing before the 5 boards see repeated bench use: either accept it (screws tightened by hand, not a lever-arm concern in practice), or note it for a respin (nudging H3 a few mm toward the bottom-right corner, or adding wire strain relief guidance in assembly notes). Visible in \texttt{bb-mcu\_top.png} and \texttt{bb-mcu\_iso.png} - compare H1/H2's tight corner hugging against H3's isolated position mid-bottom-edge.

\paragraph*{F2 (warning) - a 0.58 mm2 dead copper sliver sits on the GND layer near J3, unconnected to anything}
\texttt{kind: zero-anchor-copper-island} \textbar{} domain: \texttt{plane} \textbar{} layer B.Cu \textbar{} net GND \textbar{} refs J3 \textbar{} pos (10.158, 30.122)

The board brief for this review named this item; I did not take it on report - I re-ran `plane\_repair.py --pcb boards/bb-mcu/kicad/bb-mcu.kicad\_pcb --flag-only` myself against the built board and reproduced it independently: one GND-net component at 627.14 mm2 with 8 anchors (the real pour), and a second GND-net component at \textbf{0.5822 mm2 with 0 anchors}, centred 0.34 mm from the main pour's edge, about 0.9 mm diagonally from J3 pin 5 (GND, at 10.8, 29.48) in the board's bottom-left corner.

Electrically this is a non-event: it is DC-floating, far too small and too close to the reference pour to do anything at the frequencies on this board (SWD's own worst edge is 5 ns against 2-layer traces under 15 mm - see \texttt{architecture/blocks.md} s3), and \texttt{check\_pdn}/\texttt{check\_return\_path}/DRC all correctly see nothing wrong, because nothing electrically IS wrong. The question is fab hygiene, not function: an unconnected copper sliver 0.34 mm from its neighbour is closer to a clearance edge case than every other gap on this board, it has zero purpose (0 anchors means it grounds nothing and shields nothing), and it is exactly the sort of unexplained artifact a DFM pass or a human fab reviewer flags for manual sign-off - a turnaround-time cost for a feature that does nothing. \textbf{For the P9 DFM leg}: purge it. It has no anchors, so deleting it (or raising the zone's island-area removal threshold past 0.6 mm2 and re-filling) cannot disconnect anything real - this is a zero-risk cleanup, not a routing change requiring re-verification of anything else. Not visually distinguishable in the 3D renders at this scale (both regions are the same copper-green fill); confirmed by direct zone-fill geometry extraction, not by eye.

---

\subsubsection*{Checked and cleared (worst-first hunting, not a rubber stamp on the gates)}
{*}{*}J2 and J3 silk - verified correct and unambiguous on the render, pad by pad.{*}{*} This was the hardest item on the hunt list and I did not take \texttt{state.json}'s P6 "silk VERIFIED" decision on report. I extracted the exact absolute pad coordinates and the exact silk-text anchor coordinates from the built \texttt{.kicad\_pcb} (not estimated from pixels) and confirmed every one of the 10 net labels sits with its text-anchor at EXACTLY its pad's X coordinate (KiCad's default center-justify means this is exact centering, not approximate): J2 reads, physical pin 1 (right, square pad, silk "1" clear of the header body) to pin 5 (left): GND / SWCLK / 3V3 / SWDIO / NRST, matching the P1 ruling exactly. J3 reads pin 1 (top, square pad, silk "1" clear of the body) to pin 5 (bottom): IO1 / IO2 / IO3 / IO4 / GND, matching the P4-revised order exactly (see next paragraph). J2's two-row stagger (SWDIO/SWCLK indented below NRST/3V3/GND) is column-exact per pad, not merely close, and reads correctly at 2400 px - \texttt{bb-mcu\_top.png}, crop coordinates roughly x980-1330,y150-300 for J2 and x650-1000,y300-720 for J3 in the full-res render. Pin-1 markers (square pads plus an outside-the-body "1" digit) are present on both and would survive assembly - the exact stm32-blinky failure mode this board's own architecture doc calls out is not repeated here.

{*}{*}J2/J3 cross-plug and reversal hazards - independently re-derived, and the board matches the LATEST decision, not the one in the P4 schematic-review doc.{*}{*} \texttt{reports/review-schematic.md} W3 (as written) describes J3 as GND-at-centre and flags a rail short if a cross-plugged cable's centre wire meets J2's +3V3. That premise is now stale: \texttt{state.json} records a {*}{*}later P4 decision, triggered by that exact W3 finding{*}{*}, that moved J3's GND from centre to the end specifically to kill the rail-short case (accepting a messier reversal case - GND no longer self-maps - as the strictly safer trade, since a GPIO-to-GND/rail contact is driver-current-limited and a direct rail short is not). I re-derived the hazard tree against the AS-BUILT pinout independently (not by reading the decision text) and agree with the ruling: with J3 carrying no 3V3 pin anywhere, no reversal or cross-plug combination between these two identical unkeyed headers can short the rail to ground through the cable itself any more. \texttt{architecture/blocks.md}'s block-diagram mermaid sketch still shows the superseded IO1/IO2/GND/IO3/IO4 order - a stale diagram in an intermediate doc, not a board defect, since state.json is the decision record of authority and the built board matches it.

{*}{*}J2's one-position-slip residual risk (probe driver onto +3V3) - real, known, accepted, and its only mitigation (silk) checks out.{*}{*} \texttt{state.json} P4 already worked the alternatives and ruled there is no 5-position order safe against both a 180-reversal and a one-position slip simultaneously, and accepted the slip case as the lesser risk with per-pin silk as the sole mitigation. Silk verified above. Not re-filed as a new finding - re-litigating an already reasoned, already-accepted, correctly-mitigated risk would be noise, not signal.

{*}{*}J1 orientation and polarity silk - verified from an orthographic side render, not assumed.{*}{*} \texttt{bb-mcu\_right.png} shows the wire-entry ports as two dark rectangular openings facing the camera - i.e. facing off-board on the edge J1 actually sits on (\texttt{placement.edges} says J1 edge=right, confirmed). \texttt{bb-mcu\_left.png}/\texttt{bb-mcu\_back.png} show only the closed screw-access side from those angles, consistent with a single wire-facing side. Top-down (\texttt{bb-mcu\_top.png}, crop \textasciitilde{}x1300-1750,y150-720), "+3V3" sits fully clear above the mounted 3D terminal body and "J1 -GND" fully clear below it - the \textasciitilde{}14 mm body does not encroach on either label once assembled. Screw terminal J1 is a symmetric, unkeyed 2-pole part (confirmed in its own datasheet extraction); a 180-degree assembly rotation error is theoretically possible since nothing mechanical prevents it, but this is inherent to the part class (true of every unkeyed 2-pin THT connector, mitigated industry-wide by CAD-derived placement rotation data, not by keying) and not a defect specific to this design - considered and not filed.

\textbf{Mounting hole basics} - all four are 3.2 mm drill (M3 clearance) as required; none moved off \texttt{F.Cu}/board surface; H3's silk reference text was relocated by an earlier fix pass and reads cleanly with no new collision (\texttt{bb-mcu\_top.png}, near x1050-1150,y500-560).

\textbf{No part silently flipped to the back.} Queried every footprint's side directly (\texttt{fp.IsFlipped()}): all 14 (U1, C1-C5, R1, J1-J3, H1-H4) report \texttt{front}. Bottom render (\texttt{bb-mcu\_bottom.png}) shows only THT pin breakout, mounting-hole clearances and via witness marks on the GND pour - no SMD footprints - confirming single-sided assembly was actually delivered, not just specified.

\textbf{SWCLK does not run long-and-parallel to the four GPIO} (the one EMI note \texttt{constraints.json} flagged as advisory, unenforced by any check). Extracted every track segment on \texttt{/SWCLK}, \texttt{/SWDIO}, \texttt{/IO1..4}: SWCLK is one 4.65 mm segment confined to x=27.5, y=11.0-15.65 (top-right of the package); the four GPIO traces run entirely in x=10.8-21.5, y=18.92-26.94 (left-center). The regions do not overlap in X or Y - the architecture's own placement rationale ("maximum distance from SWCLK") holds up in the actual copper, not just in the placement note.

{*}{*}C5 (NRST cap, added at P4) is present, correctly wired, and close to its pin.{*}{*} Netlist-confirmed: C5 pad 1 on \texttt{/NRST}, pad 2 on \texttt{GND}, 3.13 mm from U1 pin 4. Resolves the schematic review's W1 (which predates C5 being added).

\textbf{Cross-artifact delivery}: all three stated interfaces present with the ruled pin orders (J1: pin1=+3V3/pin2=GND; J2: GND/SWCLK/3V3/SWDIO/NRST; J3: IO1/IO2/IO3/IO4/GND, all confirmed against the netlist, not the schematic picture). Both architecture blocks (B1 MCU minimum system, B2 SWD debug port) are present with every datasheet-required part (C1-C4, R1) on the right nets.

\textbf{One stray B.Cu track, checked and cleared}: exactly one non-zone track exists on the bottom copper layer (1.75 mm, net GND, x=29.55 near C1). It is on the SAME net as the surrounding pour, so it cannot split or interrupt anything and is not a "signal routed on the bottom" violation - almost certainly a via-stitch leftover from routing. Not filed.

---

\subsubsection*{Open / could not fully judge}
1. \textbf{F1's severity is a bench-use-frequency judgment call}, not something a render settles definitively - how hard this board actually gets handled determines whether 9.2 mm-from-nearest-hole at an open corner matters in practice. Flagged with reasoning; the call is the human's. 2. \textbf{F2's exact fab risk} (whether a 0.34 mm gap survives JLC's process without bridging, or whether their DFM software auto-flags it) needs the actual P9 DFM/fab-rules pass, not a review render - I can confirm the geometry and the recommended action (purge, zero risk) but not JLC's own tolerance behavior. 3. Did not re-render or re-read \texttt{reports/schematic.pdf} directly - relied on the existing P4 schematic review plus my own direct netlist/pad extraction from the built board (which is generated from the same schematic and is at least as authoritative for the electrical facts I needed to check). 4. Screwdriver/standoff ergonomic access at H2 given J1's \textasciitilde{}14 mm height: judged workable from the render (\textasciitilde{}2-2.5 mm clear lateral gap to the terminal body) but not physically tested.

\section{DFM and Fabrication}
Gate \texttt{dfm} (phase P9): \textbf{pass} after 2 attempt(s); last 2026-08-17T20:49:08, failing 0 of 0.
DFM findings by severity: none.
\subsection*{Fabrication outputs}
Exported layers: F.Cu, B.Cu, F.Silkscreen, B.Silkscreen, F.Mask, B.Mask, F.Paste, B.Paste, Edge.Cuts.
\subsection*{Bill of materials}
{\small
\begin{longtable}{p{4.4cm}p{2.8cm}p{4.6cm}l}
\toprule
\textbf{Comment} & \textbf{Designator} & \textbf{Footprint} & \textbf{LCSC} \\
\midrule
\endhead
100nF 50V X7R & C1,C5 & C0603 & C14663 \\
4.7uF 16V X7R & C2 & C0603 & C22399629 \\
10nF 25V X7R & C3 & C0603 & C327204 \\
1uF 16V X7R & C4 & C0603 & C106248 \\
2P 5.08mm screw terminal THT & J1 & CONN-TH\_2P-P5.00\_WJ500V-5.08-2P & C8465 \\
1x5 2.54mm male THT & J2,J3 & HDR-TH\_5P-P2.54-V-M & C32713271 \\
10k & R1 & R0603 & C25804 \\
STM32F030F4P6TR & U1 & SOP-20\_L6.5-W4.4-P0.65-LS6.4-BL & C89040 \\
\bottomrule
\end{longtable}
}
1 CPL rotation correction(s) applied; no missing LCSC numbers.
\emph{fab/order.json not available.}
\section{Run Record (Appendix)}
\subsection*{Phase digests}
\paragraph*{P0}
\begin{quote}\small\ttfamily\raggedright
\# P0 Intake - digest (2026-08-16)\\
~\\
- Mode `learning block-basics:` -\textgreater{} block-only scope, canonical binding,\\
  geometry OUTPUT (no size stated, none invented; earned at P6).\\
- `requirements.md` written; `check\_requirements.py` exit 0, 0 violations.\\
- Scope: MCU + datasheet-required support + the three connectors the brief\\
  states (J1 screw terminal 3V3 in, J2 debug, J3 4x GPIO). Protection,\\
  filtering, indicators, buttons, test points, straps, second rail: out.\\
- 3 questions asked in one batch, all defaults taken: 3.3 V is a bench/board\\
  rail and NOT a battery (last safety flag CLOSED); J2 = 5-pin\\
  3V3/SWDIO/SWCLK/NRST/GND with 3V3 sense-only; four plain digital GPIO.\\
- Owner then delegated design rulings ("best learning for the basic system"):\\
  H1-H4 become orchestrator-ruled digests. Gates, coverage, research, DFM and\\
  safety are untouched by that delegation.\\
- Fable out of credits -\textgreater{} every fable tier substituted one step up\\
  (fable/high -\textgreater{} opus/xhigh), logged in the spawn ledger.\\
- Library has 16 records, all buck-domain: expect P2/P3 coverage gaps.
\end{quote}
\paragraph*{P1}
\begin{quote}\small\ttfamily\raggedright
\# P1 Research - digest (2026-08-16)\\
~\\
- Roster: component-scout, then reference-design + interface-spec off its\\
  shortlist. power-architect skipped (one external rail = trivially powered).\\
- MCU RULED {*}{*}STM32F030F4P6TR{*}{*} (C89040, TSSOP-20, 19k stock, \$0.96). Rejected\\
  CH32V003 (SWIO needs a WCH-Link, rewrites J2), PY32 (puya.com not on the\\
  fetch allowlist), C011 (hides BOOT0 on SWCLK), G031 (68 in stock).\\
- Minimum system, cited: REQUIRED = 100 nF VDD + 4.7 uF bulk, 10 nF + 1 uF\\
  VDDA, 10k BOOT0 pull-down. RECOMMENDED-ONLY and so excluded by mode = VDDA\\
  filter, NRST cap. SWD adds zero parts. Nine parts total.\\
- SWD binds nothing on layout (5 ns edge settles \textasciitilde{}1.5 ns vs a 125 ns half-bit;\\
  JLC has no impedance-controlled 2L stackup). Guidance stays advisory.\\
- Ruled: J2 order {*}{*}GND/SWCLK/3V3/SWDIO/NRST{*}{*} (3V3-at-centre is the unique\\
  reversal-safe arrangement); no SWD series R (overruling RPi/Lauterbach/\\
  SEGGER); no NRST cap (diverges from stm32-blinky); omit the diff\_pairs key.\\
- ENV: st.com unreachable from this host (HTTP 000 after 20 s); fetch ST\\
  primaries via wmsc.lcsc.com. AN4325/RM0360/errata are on NO allowlisted\\
  host - name the host in OPEN if a record needs one. See log/env-notes.md.
\end{quote}
\paragraph*{P2}
\begin{quote}\small\ttfamily\raggedright
\# P2 Architecture + coverage research - digest (2026-08-16)\\
~\\
- Architecture written, `constraints\_lint` exit 0: 2 blocks (`mcu`,\\
  `swd-debug-port`), stackup JLC2313\_1.6 2L, one flat root sheet, 9 parts.\\
  Nothing relaxed by the mode - the brief stated no size to lose.\\
- Coverage fired as designed: first call at floor `proven` with\\
  `--research-provisional` -\textgreater{} both blocks gap, no checklist existed for either.\\
- 2 research tasks -\textgreater{} 13 records + the library's FIRST `mcu` and\\
  `swd-debug-port` checklists. 3 sources each, no forum sources, caps not hit.\\
  All 13 verified by fresh readers.\\
- The refute loop caught 3 real defects: a decoupling envelope bounded on the\\
  wrong axis (twice - now correctly `level: principle`), a false "every Min\\
  cell is blank" claim (tEXTIpw has Min 10 ns), and a family-level\\
  `external\_series\_required: false` contradicting its own sibling record.\\
- Re-run coverage: `block:B1` provisional on all 5 rows; `block:B2` gap on\\
  esd / decoupling / constraints-emission -\textgreater{} explicit designing-under-gap\\
  decision (riskless: two are classes block-only excludes).\\
- Key design finding: VDDA carries the internal RC oscillators and the PLL\\
  (DS Fig 12), so on this crystal-less board VDDA IS the clock supply and a\\
  filter into it would be a FAULT, not merely out of scope.
\end{quote}
\paragraph*{P3}
\begin{quote}\small\ttfamily\raggedright
\# P3 Parts + Library - digest (2026-08-16)\\
~\\
- `parts.json`: 9 lines (8 orderable + 1 footprint-only hole). 2 Basic /\\
  6 Extended, \textasciitilde{}1.35 USD/board. U1 stock re-verified 18,977; alt C32908.\\
- Kept ST's stated cap values though 4.7uF/10nF/1uF are all Extended at 0603\\
  (verified: JLC's Basic list carries none of them up to 1206).\\
- U1 extraction validated: 20 pins, 5 layout notes, land from Figure 37. Its\\
  pinout agrees independently with the P1 scout and P2 architect.\\
- {*}{*}DEFECT FOUND AND FIXED{*}{*}: J1's pulled footprint drilled 1.30 mm where the\\
  vendor recommends 1.50 mm; the 0.90 mm pin's square diagonal is 1.273 mm =\\
  0.027 mm clearance, a terminal that does not seat. Approved hand edit to\\
  1.50 mm drill / 2.30 mm pad. fp\_verify now 3/3 pass, 0 failed.\\
  Nothing in the pipeline would have caught it: fp\_verify had no land pattern\\
  to diff until I ordered the connector extractions, and it has no\\
  drill-vs-pin check at all.\\
- U1 land deviation ACCEPTED on worked geometry (the pulled land still\\
  captures the lead foot, with more toe and a wider inter-pad gap than ST's).\\
- Coverage: 3 part slots covered, B1 provisional, B2 gap on the 3 classes\\
  already recorded at P2. No new research opened.
\end{quote}
\paragraph*{P4}
\begin{quote}\small\ttfamily\raggedright
\# P4 Schematic - digest (2026-08-16)\\
~\\
- ONE schematic agent, one flat root sheet (recorded). `kicad/gen/root.py` is\\
  the source; the .kicad\_sch is build output.\\
- Gate {*}{*}erc PASS{*}{*} 0/0, attempt 2, commits 42f46c3 then 94d7a24.\\
  netlist\_audit exit 0: 19 nets, 10 components, 44/44 pins, 4 decouplers.\\
- The generator caught its own defect by RENDERING: J1's two labels\\
  overprinted, exporting as +3VGNID - on the connector whose silk polarity is\\
  the only defence against a swapped supply. Fixed, proven netlist-neutral.\\
- Adversarial review, fresh context: {*}{*}0 errors / 3 warnings{*}{*}, and it CLOSED\\
  the run's biggest open item - BOOT0=0 selects Main Flash, R1 is correct.\\
- W1 actioned, REVERSING my P1 ruling: C5 100 nF added on NRST. My reason had\\
  misread AN4325; the datasheet's Fig 21 note 1 recommends it unconditionally\\
  for parasitic reset, and this board flies NRST on a lead.\\
- W3 actioned: J3 -\textgreater{} IO1/IO2/IO3/IO4/GND. J2 and J3 are the SAME unkeyed part,\\
  so the GPIO cable fits the debug header; centre-GND vs centre-3V3 shorted\\
  the rail. W2 held: rail-at-an-end is worse under a one-position slip.\\
- Then caught a P6 landmine: C5 had no placement.sides entry and is SMT, so\\
  the annealer could have flipped it to the back. Added (7 now) + into the\\
  mcu group; both constraints.json copies re-synced, lint 0.
\end{quote}
\paragraph*{P5}
\begin{quote}\small\ttfamily\raggedright
\# P5 Board Setup - digest (2026-08-16)\\
~\\
- Run INLINE (no board-setup agent): no impedance work and no library repair.\\
- `board\_init --outline auto --layers 2 --stackup JLC2313\_1.6\\
  --mounting-holes 4`. Self-check: {*}{*}parity 0, setup\_violations 0{*}{*},\\
  unconnected 25 (expected, unrouted), clean true. It read the recorded build\\
  mode and honored geometry-as-OUTPUT - a fixed WxH would have been refused.\\
- Provisional outline {*}{*}51.15 x 43.524 mm{*}{*}. This is deliberately generous\\
  room, NOT a size: P6 places to the canonical layout and\\
  `board\_edit --outline fit` then earns the real dimensions.\\
- `rules\_gen` capability\_class {*}{*}2layer\_1oz{*}{*}, matching the stackup's own\\
  copper - 9 rules. Floors reported, not assumed: track and clearance\\
  0.127 mm, annular 0.150 mm, edge clearance 0.300 mm, hole-to-hole\\
  0.500 mm. J1's 0.40 mm and the header's 0.30 mm annulus both clear.\\
- Netclasses: Default only (track 0.2 mm, via 0.6/0.3). +3V3 at 0.1 A needs\\
  0.05 mm, so it buckets into Default - exactly what P2 predicted, and the\\
  reason GND was deliberately left undeclared in power{[}{]}.\\
- diff\_pairs {[}{]} and hv\_rules {[}{]} - correct, nothing on this board needs either.\\
- Sidecars beside the board: constraints.json (re-synced after the P4\\
  amendments, byte-identical to architecture/) and decoupling.json.
\end{quote}
\paragraph*{P6}
\begin{quote}\small\ttfamily\raggedright
\# P6 Placement - digest (2026-08-16)\\
~\\
- Gate {*}{*}place PASS 0/0{*}{*}, attempt 2, commit 503723c. Attempt 1 is the\\
  agent's honest FAIL against the provisional outline - see the defect below.\\
- {*}{*}PIPELINE DEFECT{*}{*}: at a geometry-OUTPUT binding the place gate must run\\
  AFTER `board\_edit --outline fit`, not before as the recipe documents.\\
  Legality measures declared edges against the CURRENT outline, so a\\
  canonically placed J1 sat 16.9 mm from the provisional edge and failed;\\
  satisfying it would spread the connectors to the provisional size, which\\
  fit then "earns" - bb-buck with the sign flipped. The fix is the ORDER.\\
- Size EARNED {*}{*}34.77 x 22.26 mm{*}{*} (774 mm2) from a 51.15 x 43.52 provisional.\\
  Hand layout beat both annealer candidates (HPWL 161.36 vs 200.12); route\\
  probe completion 1.00, 0 unrouted.\\
- U1 rot 270 is the floorplan: SWD end -\textgreater{} J2/top, VDD/VSS -\textgreater{} J1/right,\\
  PA0-PA3 -\textgreater{} J3/left, all three escapes non-crossing. Decoupler loops\\
  4.70-7.46 nH, all inside class limits.\\
- J1 rotation PROVEN by side render: its wire mouth is at footprint-local -Y,\\
  the REVERSE of the KF128 family in root LEARNINGS. Guessing would have\\
  faced the wire opening into the board - a defect no gate checks.\\
- Silk verified by me on the render: J2/J3 per-pin labels in the ruled order\\
  with pin-1 markers; J1 +3V3/-GND each beside its own pole.
\end{quote}
\paragraph*{P7}
\begin{quote}\small\ttfamily\raggedright
\# P7 Routing - digest (2026-08-16)\\
~\\
- Gate {*}{*}drc\_routed PASS 0/0{*}{*} (errors AND warnings), attempt 2,\\
  commit e01fdc1. Completion {*}{*}1.0{*}{*}, 0 unrouted nets.\\
- 2-layer chain as verified: planes\_gen -\textgreater{} route\_critical -\textgreater{} route\_auto -\textgreater{}\\
  stitch\_vias -\textgreater{} plane\_repair -\textgreater{} gate. route\_cleanup deliberately SKIPPED\\
  (2L pour board is exactly its regression class; board was already clean).\\
- planes\_gen: 1 zone, GND on B.Cu, 627 mm2. route\_critical did NOT no-op as\\
  I expected - it routed +3V3 fully, 8/8 pads at 0.3 mm via KRT.\\
- route\_auto: Freerouting ladder 3 rungs, best FR completion 0.857; KRT\\
  finish closed /NRST /SWCLK /SWDIO and GND to reach 1.0. FR's own success\\
  signal stays untrusted - only kicad-cli DRC gated it.\\
- stitch\_vias: 5 pad-stitch GND vias, one per pad, matching\\
  `mcu-decoupler-is-the-local-source`. 0 area-stitch, expected on 2L.\\
- {*}{*}The hard constraint held{*}{*}: plane\_repair splits\_found 0, the GND pour is\\
  ONE component at 627.14 mm2 with 8 anchors. Neither the anticipated\\
  SWDIO x SWCLK crossing nor the +3V3 wrap needed B.Cu at all - both\\
  resolved on F.Cu - so `mcu-supply-return-continuity` was never at risk.\\
- Carried to P8/P9: planes\_gen's fill leaves a 0.582 mm2 ZERO-ANCHOR copper\\
  island 0.34 mm from the main pour near J3's GND pad. plane\_repair correctly\\
  ruled it not a split; DRC is clean with it present. It is floating copper,\\
  so watch whether any P8 check or the P9 DFM leg flags an isolated region.
\end{quote}
\paragraph*{P8}
\begin{quote}\small\ttfamily\raggedright
\# P8 Verification - digest (2026-08-17)\\
~\\
- Gate {*}{*}verify PASS 0/0{*}{*}, attempt 2, commit da26db0. All EIGHT checks RAN\\
  and passed with NO skips - return\_path, current, decoupling, diffpair,\\
  creepage, thermal, silk, pdn. A skipped check is a hole, not a pass; there\\
  are none here. drc\_routed re-run clean (attempt 3) after the silk fix.\\
- First verify pass carried ONE warning: `silk\_misattributed` - H3's refdes\\
  sat 0.23 mm from U1 and read as U1's label. FIXED, not waived: silk\_place\\
  skips board\_only refs by design, so it needed a direct place\_edit\\
  move\_text to (29.55, 28.6), below its own hole. check\_silk now 0.\\
- Advisory legs both pass: check\_irdrop (jmax 0.63 A/mm), check\_pdn\_z.\\
- Sim gate NOT run and not applicable - no analog content, P2 recorded sim\\
  candidates as none.\\
- Board review: {*}{*}0 errors / 2 warnings / 0 waivers{*}{*}. It verified rather\\
  than trusted - re-ran plane\_repair --flag-only itself, and checked J2/J3\\
  labels by pad coordinate rather than by eye.\\
- F1 mounting-hole spread near J1: ACCEPTED (J1 braced diagonally at 9.2 mm\\
  by H2 and H3; 1.6 mm FR4 at this span). F2 copper island: carried to P9\\
  DFM with evidence rather than fixed blind.\\
- Cleared by the review: J1's wire opening faces off-board AND its polarity\\
  silk stays clear of the mounted terminal body; J2/J3 silk pixel-exact with\\
  pin 1 marked; the cross-plug rail short is genuinely gone; no silent flips.\\
- {*}{*}NO waivers exist on this board{*}{*} - verify-waivers.json was never needed.\\
- Caveat: this review ran at a FORCED tier downgrade (sonnet) - the weekly\\
  usage limit exhausted opus mid-pass. Recorded as a decision, not hidden.
\end{quote}
\paragraph*{P9}
\emph{(no digest recorded)}
\subsection*{Gate attempt history}
{\small
\begin{longtable}{llll}
\toprule
\textbf{Timestamp} & \textbf{Gate} & \textbf{Status} & \textbf{Failing/Total} \\
\midrule
\endhead
2026-08-16T21:23:29 & erc & pass & 0/0 \\
2026-08-16T21:57:21 & erc & pass & 0/0 \\
2026-08-16T22:52:31 & place & fail & 1/1 \\
2026-08-16T22:57:35 & place & pass & 0/0 \\
2026-08-16T23:15:48 & drc\_routed & pass & 0/0 \\
2026-08-16T23:16:46 & drc\_routed & pass & 0/0 \\
2026-08-16T23:17:08 & verify & pass & 0/1 \\
2026-08-16T23:20:27 & drc\_routed & pass & 0/0 \\
2026-08-16T23:20:27 & verify & pass & 0/0 \\
2026-08-17T20:48:49 & dfm & pass & 0/0 \\
2026-08-17T20:49:08 & dfm & pass & 0/0 \\
\bottomrule
\end{longtable}
}
\subsection*{Issues}
\emph{(no entries)}
\subsection*{Human checkpoints}
{\small
\begin{longtable}{lllp{2.6cm}p{6.4cm}}
\toprule
\textbf{Id} & \textbf{Phase} & \textbf{Status} & \textbf{When} & \textbf{Note} \\
\midrule
\endhead
1 & P2 & approved & 2026-08-16T20:18:19 & DELEGATED RULING, not an owner approval: the owner delegated design rulings for this run ('make decisions yourself based on what will give the best learning for the basic system'). Orchestrator ruled H1 pass: 2 blocks, JLC2313\_1.6 2L, 9 parts, \textasciitilde{}1.30 USD/board BOM, nothing relaxed by the mode (no size was ever stated). Coverage: block:B1 provisional on all 5 rows, block:B2 gap on esd/decoupling/constraints-emission - recorded as an explicit designing-under-coverage-gap decision. Research caps NOT hit: 4 of 6 tasks and 1 of 4 depth units per task left unspent. Riskiest decision: omitting the NRST 100 nF cap, diverging from the shipped stm32-blinky. Owner items surfaced but not blocking: the fetch allowlist (st.com and analog.com are unreachable from this host though allowlisted; arm.com and segger.com are absent and would be needed to source the probe side), and the promotion ruling on 13 verified records + 2 new checklists \\
2 & P4 & approved & 2026-08-16T21:58:16 & DELEGATED RULING, not an owner approval. Orchestrator ruled H2 pass. erc gate PASS (0/0, attempt 2, commit 94d7a24), netlist\_audit exit 0, 10 parts, 44/44 pins. Reviewer 0 errors / 3 warnings, all three ruled: W1 NRST cap ADDED (reversing my own P1 ruling - the reviewer showed my reason misread AN4325 where the datasheet Fig 21 note 1 recommends it unconditionally for parasitic reset, and this board flies NRST on a lead); W3 J3 reordered to move GND off centre (J2/J3 are the same unkeyed part, so the GPIO cable cross-plugged into the debug header shorted the rail); W2 J2 order held unchanged with reasoning recorded (rail-at-an-end is strictly worse under a one-position slip). Review also CLOSED the run's BOOT0 polarity open item: BOOT0=0 selects Main Flash, the 10k pull-down is correct. No ERC waivers exist - there are no ERC warnings to waive. Carried to P6 as a hard requirement: per-pin silk on J2 and J3, which is the only mitigation for the unkeyed-header residual risk \\
3 & P6 & approved & 2026-08-16T22:59:47 & DELEGATED RULING, not an owner approval. H3 (optional, default ON) ruled pass on the fitted top render, inspected directly. place gate PASS 0/0 attempt 2, commit 503723c. Board EARNED 34.77 x 22.26 mm from a 51.15 x 43.52 provisional - 65 percent area reduction, and the size is an output of the layout as the binding requires. Route probe 1.00 with 0 unrouted. Silk checked on the render not taken on report: J2 and J3 per-pin labels both in the ruled order with pin-1 markers, J1 polarity legible with each label beside its own pole. Accepted the agent's overrule on moving H1-H4 (leaving them at provisional corners would have defeated the fit) and its rejection of four true corner holes (+25 percent area of pure air on a board with no enclosure requirement). Residual accepted: rightmost \textasciitilde{}10 mm of the bottom edge unsupported, J1 braced by H2 above and H3 below-left. Found and recorded a PIPELINE DEFECT in the recipe's gate ordering at geometry-OUTPUT bindings \\
4 & P8 & approved & 2026-08-17T20:47:52 & DELEGATED RULING, not an owner approval. H4 ruled pass. Machine state fully green with NO skips: erc 0/0, place 0/0, drc\_routed 0/0, verify 0/0 across all eight checks, check\_irdrop pass (jmax 0.63 A/mm), check\_pdn\_z pass. Board review 0 errors / 2 warnings / 0 waivers recommended - both warnings ruled: F1 mounting-hole spread accepted (J1 braced diagonally at 9.2 mm by H2 and H3), F2 copper island carried to P9 DFM with evidence rather than fixed blind. Reviewer independently cleared the items that decide bench usability: J1's wire opening faces off-board and its polarity silk stays clear of the mounted terminal body; J2/J3 per-pin silk is pixel-exact on its pads with pin 1 marked; the cross-plug rail-short is genuinely resolved by the P4 J3 reorder; no part silently flipped to back; C5 present 3.1 mm from U1 pin 4. NO waivers exist on this board - reports/verify-waivers.json was never needed. Caveat recorded separately: this review ran at a FORCED tier downgrade (sonnet) because the weekly usage limit exhausted opus mid-pass \\
5 & P10 & not reached &  &  \\
\bottomrule
\end{longtable}
}
\section{Artifact Index}
{\small
\begin{longtable}{p{9cm}rp{4cm}}
\toprule
\textbf{File} & \textbf{Size} & \textbf{Note} \\
\midrule
\endhead
\texttt{state.json} & 76,436 B &  \\
\texttt{kicad/bb-mcu.kicad\_pcb} & 134,059 B &  \\
\texttt{kicad/bb-mcu.kicad\_sch} & 60,588 B &  \\
\texttt{fab/bb-mcu\_gerbers.zip} & 20,201 B &  \\
\texttt{fab/BOM.csv} & 376 B &  \\
\texttt{fab/CPL.csv} & 379 B &  \\
\texttt{brief/brief.md} & 374 B &  \\
\texttt{requirements.md} & 19,063 B &  \\
\texttt{architecture/blocks.md} & 19,238 B &  \\
\texttt{architecture/constraints.json} & 16,927 B &  \\
\texttt{architecture/power\_tree.md} & 10,279 B &  \\
\texttt{architecture/sheets.md} & 12,412 B &  \\
\texttt{architecture/stackup.md} & 8,344 B &  \\
\texttt{reports/board\_edit-fit.json} & 2,016 B &  \\
\texttt{reports/board\_edit-probe.json} & 1,965 B &  \\
\texttt{reports/bom\_cpl.json} & 7,106 B &  \\
\texttt{reports/check\_irdrop.json} & 2,085 B &  \\
\texttt{reports/check\_pdn\_z.json} & 893 B &  \\
\texttt{reports/check\_silk.json} & 537 B &  \\
\texttt{reports/fab\_export.json} & 2,216 B &  \\
\texttt{reports/gate-dfm.json} & 14,935 B &  \\
\texttt{reports/gate-drc\_routed.json} & 1,072 B &  \\
\texttt{reports/gate-erc.json} & 14,213 B &  \\
\texttt{reports/gate-place.json} & 9,605 B &  \\
\texttt{reports/gate-verify.json} & 11,175 B &  \\
\texttt{reports/render-fitted.json} & 447 B &  \\
\texttt{reports/render-verify.json} & 451 B &  \\
\texttt{reports/review-board.json} & 1,927 B &  \\
\texttt{reports/review-board.md} & 13,039 B &  \\
\texttt{reports/review-schematic.json} & 2,883 B &  \\
\texttt{reports/review-schematic.md} & 20,832 B &  \\
\texttt{reports/rules\_gen.json} & 804 B &  \\
\texttt{reports/schematic.pdf} & 62,739 B &  \\
\texttt{reports/silk-fix-holes.json} & 946 B &  \\
\texttt{reports/silk-h3-fix-report.json} & 255 B &  \\
\texttt{reports/silk-h3-fix.ops.json} & 124 B &  \\
\texttt{reports/top.net} & 24,807 B &  \\
\bottomrule
\end{longtable}
}
\end{document}
